\documentclass[11pt,twoside]{article}
\usepackage{latexsym}
\usepackage{amssymb,amsbsy,amsmath,amsfonts,amssymb,amscd}
\usepackage{dsfont}
\usepackage{graphicx,color}
\usepackage[colorlinks=true]{hyperref}
\usepackage{epsfig,color}
\usepackage{dsfont}
\usepackage{algorithm}
%\usepackage{algorithm2e}
%\usepackage{algoritmic}
%\usepackage{algpseudocode}

\usepackage[amsmath,framed,thmmarks]{ntheorem}
\theoremstyle{nonumberplain}


\usepackage{subcaption}
\usepackage{amsmath} 
\DeclareMathOperator{\sign}{sign}
\usepackage{longtable, array, booktabs}
\newcommand\red[1]{\color{red}#1} 
\usepackage{tabularx}
\usepackage{mathtools}
\setlength{\oddsidemargin}{0mm}
\setlength{\evensidemargin}{0mm}
\setlength{\topmargin}{5mm}
\setlength{\textheight}{20cm}
\setlength{\textwidth}{15cm}
\renewcommand{\baselinestretch}{1.0}
\parindent 10pt

\DeclareMathOperator{\Tr}{Tr}

\DeclareMathOperator*{\argmax}{arg\,max}
\DeclareMathOperator*{\argmin}{arg\,min}

\newcommand*\diff{\mathop{}\!\mathrm{d}}

\newtheorem{theorem}{Theorem}%[section]
\newtheorem{lemma}[theorem]{Lemma}

\newtheorem{proof}[theorem]{Proof}


\newtheorem{definition}[theorem]{Definition}
\newtheorem{remark}[theorem]{Remark}
\newtheorem{proposition}[theorem]{Proposition}
\newtheorem{corollary}[theorem]{Corollary}



\newcommand{\bd}{\mbf{d}}
\newcommand{\bu}{\mbf{u}}
\newcommand{\bv}{{\mathbf{v}}}
\newcommand{\bw}{\mbf{w}}
\newcommand{\bx}{{\mathbf{x}}}
\newcommand{\bX}{\mbf{X}}
\newcommand{\bbX}{\mathbb{X}}
\newcommand{\bY}{\mbf{Y}}
\newcommand{\mH}{\mathcal{H}}

\newcommand{\rhoT}{\rho_T}
\newcommand{\rhoL}{\rho_T^L}
\newcommand{\rhoLM}{\rho_T^{L,M}}
\newcommand{\brhoL}{\bm{\rho}_T^L}
\newcommand{\rhsF}{\mbf{f}}

\newcommand{\brho}{\boldsymbol{\rho}}
% customized commands


\newcommand{\mres}{\mathbin{\vrule height 1.2ex depth 0pt width
0.13ex\vrule height 0.13ex depth 0pt width 0.9ex}}

\definecolor{myblue}{rgb}{0.0352,0.4981,0.6509}
\newcommand{\GER}{\color{myblue}}

\newcommand{\mbf}[1]{\boldsymbol{#1}}
\newcommand{\lil}[1]{\ell_{#1}}
\newcommand{\innerp}[2]{\langle #1,#2 \rangle}
\newcommand{\inp}[1]{\langle{#1}\rangle}
\newcommand{\dbinnerp}[1]{\langle\hspace{-1mm}\langle{#1}\rangle\hspace{-1mm}\rangle}
\newcommand{\abs}[1]{\big| #1 \big|}
\newcommand{\dquo}[1]{``#1"}
\newcommand{\bxm}{\bx^{m}}
\newcommand{\dotbxm}{\dot\bx^m}
\newcommand{\bXm}{\bX^{m}}
\newcommand{\dotbXm}{\dot\bX^{m}}
\newcommand{\Xtrain}{\bX_{\mathrm{tr}}}
\newcommand{\brm}{\bar{\br}^{m}}
\newcommand{\nbrm}{r^{m}}

% define norms
\newcommand{\norm}[1]{\left\| #1 \right\|}
\newcommand{\Tnorm}[1]{\left\| #1 \right\|_{\mathcal{T}}}
\newcommand{\tnorm}[1]{\left\| #1 \right\|_{\mathcal{T}_t}}
\newcommand{\HDnorm}[1]{d_{\text{HD}}{(#1)}}
\newcommand{\ltworho}[1]{\norm{#1}_{L^2(\rho_T)}}
% define some spaces in real
\newcommand{\realR}[1]{\mathbb{R}^{#1}}
\newcommand{\real}{\mathbb{R}}
\newcommand{\compC}[1]{\mathbb{C}^{#1}}
\newcommand{\comp}{\mathbb{C}}
\newcommand{\integ}{\mathbb{Z}}
% for some special variables
% the vector varaible


\newcommand{\br}{\mbf{r}}

\newcommand{\bB}{\mbf{B}}

%\newcommand{\bX}{\mbf{X}}



\newcommand{\balpha}{\mbf{\alpha}}
%
\newcommand{\mA}{\mathcal{A}}
\newcommand{\mC}{\mathcal{C}}
\newcommand{\mE}{\mathcal{E}}
\newcommand{\mF}{\mathcal{F}}
\newcommand{\mG}{\mathcal{G}}
\newcommand{\mL}{\mathcal{L}}
%\newcommand{\mH}{\mathcal{H}}

\newcommand{\mI}{\mathcal{I}}
\newcommand{\mK}{\mathcal{K}}
\newcommand{\mR}{\mathcal{R}}
\newcommand{\mS}{\mathcal{S}}
\newcommand{\mT}{\mathcal{T}}
\newcommand{\mX}{\mathcal{X}}
\newcommand{\mY}{\mathcal{Y}}
%
\newcommand{\bigO}{\mathcal{O}}
%
\newcommand{\R}{\real}
\newcommand{\dimamb}{D}
\renewcommand{\dim}{d}
\newcommand{\numcl}{K}
\newcommand{\idxcl}{k}
\newcommand{\spaceM}{S}
\newcommand{\cl}{C}
\newcommand{\clof}{\mathpzc{k}}
\newcommand{\adj}{\mathcal{A}}
\newcommand{\eadj}{\adj^E}
\newcommand{\aadj}{\adj^A}
\newcommand{\xiadj}{\adj^{\xi}}
% the coupling strength terms
\newcommand{\cst}{\mathcal{\kappa}}
\newcommand{\vcst}{\cst^{\bv}}
\newcommand{\xicst}{\cst^{\xi}}
\newcommand{\bcst}{\bm{\cst}}
\newcommand{\bvcst}{\bm{\cst}^{\bv}}
\newcommand{\bxicst}{\bm{\cst}^{\xi}}
\newcommand{\bL}{\mbf{L}}
\newcommand{\e}{\varepsilon}

%

\newcommand{\st}[1]{\textcolor{red}{#1}}
\newcommand{\XR}[1]{\textcolor{blue}{#1}}



\title{An asymptotic preserving scheme for kinetic equations describing bacterial travelling pulses}
\author{Discussion notes}
\date{August 2021}

\begin{document}

\maketitle



\section{Introduction}


Chemotaxis is the directed movement of bacteria in response to chemical fluctuations in the environment. Many bacteria navigate towards favourable conditions by sensing gradients of nutrients or chemoattractants and away from harmful substances. This behaviour is essential for survival \cite{franca2006navigating} and it plays a crucial role in many biological processes such as wound healing, fertilization and inflammation \cite{rickert2000leukocytes}, in multicellular organisms, and for unicellular organisms is used as a foraging mechanism \cite{willard2006signaling} \cite{franca2010g}.



Although bacterial chemotaxis has been studied since the end of the 19th century,  it was not until the work of E.F. Keller and L.A. Segel in 1970 that a precise mathematical model was proposed \cite{keller1971model}. This result became a cornerstone in the field of mathematical biology and remains one of the most influential and widely cited to this day, thanks to the rigorous coupled partial differential equation (PDE) system introduced for modelling chemotaxis.

The general form of Keller-Segel (K-S) model is
\begin{equation}
      \left \{
    \begin{aligned}
    \frac{\partial \rho}{\partial t} &= \nabla (D_\rho (\rho, S) \nabla \rho - \chi (\rho, S) \nabla S) + f(\rho, S), \\
    \frac{\partial S}{\partial t} &= D_S \Delta S + g (\rho, S) - h(\rho, S)S.
    \end{aligned}
  \right .
  \label{ks_model}
\end{equation}

Here, $\rho (t, x)$ represents the bacteria density on a given domain $\Omega \subset \mathbb{R}^n $ and $S(t, x)$ denotes the concentration of the chemical signal. The first term in the first equation in (\ref{ks_model}) describes the random component of the movement, while the second term describes the bias component, bacteria move towards high concentrations of the chemical. Consequently, the function $D_\rho (\rho, S)$ describes the diffusivity of the bacteria (random component), and $\chi (\rho,S)$ represents the chemotactic sensitivity, which measures how sensible bacteria are to a chemical variation (biased component). Both $D_\rho$ and $\chi$ may be functions of $\rho$ and $S$. The function $f(\rho, S)$ accounts for bacteria growth and death, whereas the functions $g(\rho, S)$ and $h(\rho, S)$ are functions that describe production and degradation of the chemical signal, respectively. Finally, $D_S$ is a diffusion coefficient of $S$.


We study the migration of \textit{E.~coli} driven by high concentrations of nutrient $N$ and chemoattractant $S$, motivated by the microchannel experiment of \cite{saragosti2010mathematical}. Roughly $10^5$ cells are initially concentrated at the left end of a $2\,\mathrm{cm}$ channel (width $500\,\mu\mathrm{m}$, height $100\,\mu\mathrm{m}$); nutrient is supplied at the right end. On the $\sim 100\,\mathrm{s}$ time scale, a dense pulse forms and propagates rightward with nearly constant speed and mass. The profile is strongly asymmetric—diffuse at the front and sharp at the rear—and this shape is preserved during propagation. Cell division is neglected because the experimental duration (hours) is short relative to the cell cycle.


We numerically investigate whether the relevant macroscopic system of PDEs corresponds to the underlying physical system described by the kinetic equation {in the diffusion limit} and to verify the emergence of travelling wave solutions at both the kinetic and macroscopic levels. These methods are widely known as asymptotic preserving schemes (AP) \cite{filbet2010class,jin1999efficient} since they mimic the asymptotic behaviour of the kinetic equation when the scaling parameter approaches to zero and the mesh size and time steps are fixed. In this paper, we use a micro-macro decomposition of the unknown in the sense of  \cite{lemou2008new} as detailed in Section ??.\\




Following the seminal works by Alt \cite{alt1980biased}, Dunbar and Othmer \cite{othmer1988models}  we start from the following kinetic equation 
\begin{equation}
    \partial_t f(t,x,v)+{v}\partial_x f(t,x,v)= \int_{V}\Bigl(T(t,x,v,v')f(t,x,v')-  T(t,x,{v}',v)f(t,x,v)\Bigr)\diff v'\ .\label{eq: main kinetic}
\end{equation}
The bacteria population is described by the mesoscopic density $f(t,x,v)$ where $(x,v)\in\mathbb{R}\times V$ since we are looking for planar travelling wave solutions. Here $V$ denotes a compact interval of admissible velocities. The turning kernel $T$ describes the change in direction during the tumbling phase of the movement. $T$ depends on the space and time variables through the chemical signals in the medium,  and $T(\cdot, v,v')$ represents the change of velocity from $v'$ to $v$. Following [Refs], we consider the influence of two chemical species, $S(t,x)$, a chemotactic signal released by the bacteria, and $N(t,x)$, a nutrient consumed by the bacteria. We assume that both chemical cues influence the tumbling kernel in an additive way as we explain below. For simplicity we are going to assume that $T$ only depends on the current velocity with the condition that $\int_VT(t,x,v)\diff v=1$.

Let 
\begin{equation}\label{TepsFeps}
%\begin{cases}
    T_\varepsilon(t,x,v)=\frac{\psi_\varepsilon(t,x,{v}) }{\displaystyle\int_{V}\psi_\varepsilon(t,x,{v})  \diff {v}}\ ,
 %   \end{cases}
\end{equation}
where we assume that the dependence of the turning operator on the chemical concentrations $S$ and nutrients $N$ happens as a perturbation of magnitude $\varepsilon\ll 1$ in the following way
\begin{equation}\label{expansion_psi}
    \psi_\varepsilon(t,x,{v})=1+\varepsilon\psi_1(t,x,{v})\ ,
\end{equation}
%\paragraph{Assumption 1}
where 
\begin{equation}
\psi_1(t,x,{v}) = \phi_1[S]({v}) + \phi_2[N]({v})\ , \quad\int_V\psi_1(t,x,{v})\diff {v}=0 \ , \label{hypo2}
\end{equation}
and $\phi_1$ and $\phi_2$ are given by
$$\phi_1[S]({v})=\chi_S\phi\Bigl(\frac{DS}{Dt}\bigg\rvert_{v}\Bigr)\ \ ,\ \ \phi_2[N](v)= \chi_N\phi\Bigl(\frac{DN}{Dt}\bigg\rvert_{v}\Bigr)\ .$$
 Here $\frac{D}{Dt}=\partial_t+{v}\partial_x$ is the total derivative along the direction given by ${v}$.


 
Using \eqref{expansion_psi}, the turning kernel $T_\varepsilon$ writes
\begin{align}
    T_\varepsilon(t,x,{v})=& 1+\varepsilon\Bigl[\chi_S\phi\Bigl(\frac{DS}{Dt}\bigg\rvert_{v}\Bigr)+ \chi_N\phi\Bigl(\frac{DN}{Dt}\bigg\rvert_{v}\Bigr)\Bigr]  + {\mathcal{O}(\varepsilon^2)} \ ,\label{eq: turning operator scaled}
\end{align}
{where the error terms contained in $\mathcal{O}(\varepsilon^2)$ integrate to $0$ in the velocity variable. We also note that using \eqref{hypo2} we have $\int_V T_\varepsilon(t,x,v)\diff v=1$ which describes the conservation of individuals during the velocity reorientation.}

The chemoattractant and nutrients concentrations evolve according to a reaction-diffusion equation
\begin{align}
    \partial_t S&=D_S\partial_{xx}S-\alpha S+\beta \rho\ ,\\
    \partial_t N&=D_N\partial_{xx}N-\gamma\rho N\ .
\end{align}
The chemoattractant $S$ is consumed and produced by the cells with rates given by $\alpha$ and $\beta$, respectively. The nutrients are consumed at rate $\gamma$. Both, chemoattractant and nutrients diffuse in the medium with constant $D_S$ and $D_N$. Here $\rho$ denotes the macroscopic density defined as
\[
\rho(t,x)=\int_Vf(t,x,v)\diff v\ .
\]

As we mention before, we are going to study this biological phenomenon using an approach distinct from classical modelling. In particular, the biased component in the equation of $\rho$ in (\ref{ks_model}), which is going to be derived from a kinetic description of the motion, will be uniformly bounded, whereas in the case of K-S this biased component is the chemotactic sensitivity function $\chi(\rho, S)$, which is proportional to the gradient of the chemical, is generally unbounded when instability due to aggregation takes place. This will provide us with a more precise and rigorous description of the evolution of baterial movement in this specific context, what is to say, it will allow us to describe analytically the travelling wave behaviour it is observed experimentally.

In this paper we numerically study the macroscopic limit of the kinetic equation \eqref{eq: main kinetic} to the macroscopic PDE
\begin{equation}
  \left \{
    \begin{aligned}
       \partial_t \rho &= D_\rho \Delta \rho - \nabla \cdot (\rho u_S + \rho u_N),\\
       \partial_t S &= D_S \Delta S - \alpha S + \beta \rho,\\
        \partial_t N &= D_N \Delta N - \gamma \rho N.
    \end{aligned}
  \right .
  \label{basicmodel}
\end{equation}

On the one hand, the density of bacteria $\rho(t, x)$ (at time $t>0$ and position $x \in \mathbb{R} $) follows a drift-diffusion equation, combining linear diffusion (first of the RHS) together with transport processes (second term in the RHS). On the other hand, we consider the chemoattractant signal $S(t,x)$ and the nutrient $N(t,x)$ which both follow reaction-diffusion equations driving the external chemical concentrations. The latter term of the first equation includes $u_S$ and $u_N$ that are fluxes coming from the chemoattractant and the nutrient, respectively, and which are responsible for the biased movement of the bacteria. In particular, the flux $u_S$ contributes to gather the bacterial density and create a pulse, and the flux $u_N$ is responsible for the motion of this pulse towards higher nutrient levels.
Finally, we consider that the chemoattractant is naturally degraded at rate $\alpha$ and produced by the bacteria at rate $\beta$. Also, the nutrients are consumed at rate $\gamma$. 

In contrast to the definition of the fluxes in the K-S system, given by

\begin{equation}
    u_S = \chi(\rho, S) \nabla S, \quad \quad u_N = \chi(\rho, N) \nabla N,
    \nonumber
\end{equation}

here we propose to follow the analysis in \cite{dolak2005kinetic}, in which the fluxes are 
\begin{equation} \label{fluxes}
\begin{split}
     u_S &= \chi_S J_\phi (-\epsilon \partial_t S, |\nabla S|) \frac{\nabla S}{|\nabla S|}, \\ 
    u_N &= \chi_N J_\phi (-\epsilon \partial_t N, |\nabla N|) \frac{\nabla N}{|\nabla N|},   
\end{split}
\end{equation}
based on the temporal response of bacteria $\phi$, we refer to Section \ref{sec: maro system} for a detailed definition of this function. This $\phi$ is called a response function, and describes how a single bacteria responds to a change in the concentration of the chemoattractant $S$ (or of the nutrient $N$) in the surrounding environment. The function $J_\phi$ incorporates the microscopic features that derive from the concrete response of each bacterium $\phi$ to a change in the environment. The parameter $\epsilon$ is the ratio between the pulse speed and the speed of
individual bacteria, in the experiments we can see that $\epsilon \approx 0.1$, which means that the pulse travels one order of magnitude slower than each individual. Moreover, the constants $\chi_S$ and $ \chi_N$ are the chemotactic biases, which characterize the amplitude of the signal integration function $\phi$. Finally, it is worth noting that the fluxes $u_S$ and $u_N$ are uniformly bounded, as we were expecting.

This paper is organised as follows. \textcolor{red}{(...)}


\section{Macroscopic system of equations}\label{sec: maro system}
In this section we are going to derive the macroscopic system that is obtained from \eqref{eq: main kinetic} following \cite{saragosti2010mathematical}. We introduce two relevant parameters: the maximum speed of a single bacterium $c = \max\{|v|: v \in V\}$ and the mean turning frequency $\lambda_0$.

The main unknown is the speed of the travelling pulse, $\sigma$. We start by rescaling the kinetic model (\ref{eq: main kinetic}) into a non-dimensional form using the non-dimensional variables, denoted with tildes, as
\begin{equation}
    \tilde{t}= \frac{{t}}{\bar{t}}\ , \quad \tilde{x} = \frac{x}{\bar{x}}\ , \quad \tilde{v} = \frac{v}{c}\ , \quad \tilde{T}=\frac{T}{\lambda_0}\ ,
    \nonumber
\end{equation}
to model the travelling wave with $\bar{x} = \sigma\bar{t}$. According to the experiments the pulse
speed $\sigma$ is much lower than the speed of a single bacterium $c$, and this difference is measured by the ratio $\epsilon = \sigma / c$, where $\epsilon \ll 1$.



The kinetic equation can be written in the following form, where we  drop the tilde for simplicity of notation 
\begin{equation}
    \varepsilon \partial_t f^\varepsilon+{v}\partial_x f^\varepsilon = \frac{\mu}{\varepsilon}\int_V\Bigl(T_\varepsilon(t,x,v')f^\varepsilon(t,x,v')-T_\varepsilon(t,x,v)f^\varepsilon(t,x,v)\Bigr)\diff v'\ , \label{eq: inital model with scaling}
\end{equation}
where $\mu/\varepsilon=\lambda_0 \bar{x}/c$.
We now consider a second order regular expansion of $f^\e$ in $\varepsilon$, 
\begin{equation*}
f^\e(t,{x},{v})=f^0(t,{x},{v})+\varepsilon f^1(t,{x},{v})+\varepsilon^2 f^2(t,{x},{v})+\mathcal{O}(\varepsilon^3)\ , 
\end{equation*}
where $\int_V f^\e(t,{x},{v})\diff {v} = \int_V f^0(t,{x},{v})\diff {v} = \rho(t,{x})$, therefore $\int_V f^i(t,{x},{v})\diff {v}=0$, $\forall i\geq 1$.






Introducing this ansatz in \eqref{eq: inital model with scaling}, we obtain
\iffalse
{\begin{align*}
    \varepsilon^2\partial_t f^0&+\varepsilon{v}\partial_x\big(f^0+\varepsilon f^1  \big) 
    =-\psi_0(v) f^0+\int_V\psi_0(v')f^0(v')\diff v'\\
    &+\varepsilon\int_V\Bigl(\psi_0(v')f^1(v')+\psi_1(t,x,v')f^0(v')-\psi_0(v)f^1(v)-\psi_1(t,x,v)f^0(v) \Bigr)\diff v'+{\mathcal{O}(\varepsilon^3)}
      \ .
\end{align*}}

\noindent Identifying the different equations in powers of $\varepsilon$, we obtain
\fi
\begin{align}
    \varepsilon^0: & \quad f^0(t,{x},{v})=\rho(t,{x})\ ,\label{eq: epsilon0}\\
    \varepsilon^1: & \quad {v}\partial_x f^0=\int_V\psi_1(v')f^0(v')\diff v'{+}\psi_1(v)f^0(v)-f^1(v)\ ,\label{eq: epsilon1}\\
    \varepsilon^2: & \quad \partial_tf^0+{v}\partial_x f^1 =\int_V\psi_1(v')f^1(v')\diff v'-\psi_1(v)f^1(v)-f^2(v) \ \label{eq: epsilon2}.\
\end{align}
Next, after replacing $f^0$ by its expression (Eq. \eqref{eq: epsilon0}), we multiply \eqref{eq: epsilon1} by ${v}$ and we integrate again with respect to ${v}$ to obtain
\begin{align}
    \int_V vf^1(t,x,v) \diff v=&-\partial_x\rho(t,x)\int_V |v|^2\diff v-\rho(t,x)\chi_S\int_V\phi(\varepsilon\partial_tS+v\partial_xS)v\diff v\nonumber\\ &-\rho(t,x)\chi_N\int_V\phi(\varepsilon\partial_tN+v\partial_xN)v\diff v.\label{eq: mean direction}
\end{align}
We define
\begin{equation}
D_\rho=\int_V v^2 \diff v\ ,\quad u_S=\chi_S\int_V\phi(\varepsilon\partial_tS+v\partial_xS)v\diff v \ ,\quad \textnormal{and}\quad u_N=\chi_N\int_V\phi(\varepsilon\partial_tN+v\partial_xN)v\diff v\ .\label{eq: moment}\end{equation} 
In the limiting case where the internal response function $\phi$ is bivaluated, specifically defined as
\[
\phi(Y) = \phi_0 \mathbf{1}_{\{Y < 0\}} - \phi_0 \mathbf{1}_{\{Y > 0\}} = - \phi_0 \, \sign(Y),
\]
the flux terms simplify to
\begin{equation}
u_S= \chi_S\sign(\partial_x S)\ , \quad u_N=\chi_N \sign(\partial_x N)\ .
\end{equation}
Finally, integrating \eqref{eq: epsilon2} in ${v}$ and using \eqref{eq: moment} we obtain
\[
\partial_t\rho=D_\rho\partial_{xx}\rho-\partial_x(\rho u_S+\rho u_N)\ .
\]
We refer to \cite{emako2015existence,emako2016traveling} for more details of the derivation.


\iffalse
Now we non-dimensionalised the system. We introduce a new set of variables $\tilde{t}=t/\bar{t}$ and $\tilde{x}=x/\bar{x}$ where $\bar{t}$ and $\bar{x}$ are the macroscopic time and space scales. Then we have
\begin{align}
&\partial_t\rho =\tilde{D}_\rho\partial_{xx}\rho-\partial_x (\rho \tilde{u}_S+\rho \tilde{u}_N)\ ,\label{eq: rho nondimension}\\
&\partial_tS =\tilde{D}_S\partial_{xx} S-\tilde{\alpha} S+\tilde{\beta} \rho\ ,\\
&\partial_t N=\tilde{D}_N\partial_{xx} N-\tilde{\gamma} \rho N\ \label{eq: N nondimension},
\end{align}
where $\tilde{u}_S=\frac{\bar{t}\chi_S}{\bar{x}}\sign(\partial_x S)$, $\tilde{u}_N=\frac{\bar{t}\chi_N}{\bar{x}}\sign(\partial_x N)$ and 
\[
\tilde{D}_\rho=\frac{\bar{t}D_\rho}{\bar{x}^2}\ ,\ \ \tilde{D_S}=\frac{\bar{t}D_S}{\bar{x}^2}\ ,\ \ \tilde{\alpha}=\bar{t}\alpha\ ,\ \ \tilde{\beta}=\bar{t}\beta\ ,\ \ \tilde{D}_N=\frac{\bar{t}D_N}{\bar{x}^2}\ ,\ \ \tilde{\gamma}=\gamma \bar{t}\ .
\]
From now on, for simplicity of notation we are going to drop the tilde in the new coefficients.
\fi

%\section{Travelling pulses under competing fluxes}

%\textcolor{red}{(To Canan: In this section we should include the traveling wave solution for the kinetic and macroscopic systems. Please, we should write things clearly and not copy and paste results from previous papers. Since everything is already done in \cite{saragosti2010mathematical} for the macroscopic equations and in \cite{calvez2018well} for the kinetic equation we should only include here the main results: traveling wave solution, speed of propagation etc)}


\section{Existence of travelling wave solutions}

In this section we summarise the results obtained in \cite{calvez2016chemotactic} for the existence of travelling wave solutions for the kinetic equation \eqref{eq: main kinetic},  and for the macroscopic system \eqref{basicmodel}.

Consider a travelling wave solution $z=x-ct$ where $c$ denotes the speed of the wave. We rewrite \eqref{eq: main kinetic} in the moving frame variable $z$,
\begin{equation}
\begin{aligned}\label{system}
(v-c)\partial_zf(z,v)&=\int T(z,v'-c)f(z,v')d\nu(v')-T(z,v-c)f(z,v)\\
-c\partial_z S(z)&=D_S\partial_{z}^2S(z)-\alpha S(z)+\beta \rho(z)\\
-c\partial_zN(z)&=D_N\partial_{z}^2N(z)-\gamma\rho(z) N(z)
,
\end{aligned}
\end{equation}
where $\nu$ is a symmetric probability measure on $\mathbb{R}^d$, where $V=\textnormal{supp}\ \nu$. Throughout this section we assume that $\nu$ is absolutely continuous with respect to the Lebesgue measure. The probability density function belongs to $L^p$ for some $p > 1$, i.e.
\[
\diff \nu (v)=w(v)\diff v\ ,\ \quad w\in L^p\ , \quad p\in(1,\infty]\ .
\]

We assume that the the fresh nutrient is located on the right so the wave moves towards high concentrations of the nutrient, i.e.  $\partial_zN>0$. We expect the chemoattractant concentration to have a maximum at the cell density peak, and therefore we look for a solution where  $\partial_zS$ changes sign. Without loss of generality we assume that this maximal point is at $z=0$.  We summarise these conditions as follows
\begin{equation}
\begin{aligned}
 &S(\pm \infty)=0\ ,\quad N(\pm\infty)=N_\pm\ ,\\
&\partial_z S>0  \quad  \textrm{for}\quad z<0\ , \quad\textnormal{and}\quad\partial_z S<0 \quad \textrm{for}\quad z>0\ ,\\
&\partial_zN>0\ ,\quad\forall z\in\mathbb{R}\ .\label{eq: cond chemicals}
\end{aligned}
\end{equation}
Moreover, the tumbling rate $T$ in \eqref{system} can take only four values depending on the sign of the gradients,
\begin{eqnarray}\label{tumbl}
   T(z,v-c)=1-\chi_S \textrm{sign}((v-c)\partial_z S(z))-\chi_N \textrm{sign}((v-c)\partial_z N(z))\ ,
\end{eqnarray}

\paragraph{Macroscopic travelling pulse}
We present a mathematical description of travelling pulses at the macroscopic scale. For the proof of this theorem we refer the reader to \cite{calvez2016chemotactic}.

\begin{theorem}[\cite{calvez2016chemotactic}]\label{thmacro}
There exist a speed $\sigma>0$, and two positive values $N_- < N_+$, such that the system (\ref{basicmodel}) admits travelling wave solutions ($\rho (x- \sigma t), S(x - \sigma t), N(x - \sigma t)$) with the following boundary conditions:
\begin{equation}
    \lim_{|z| \to \infty} \rho = \lim_{|z| \to \infty} S = 0, \hspace{0.6cm} \lim_{z \to -\infty} N = N_- < N_+ = \lim_{z \to + \infty} N.
    \nonumber
\end{equation}
Assume in addition that $\partial_z S$ changes sign only once, and that $\partial_z N$ is positive everywhere. Then, the speed $\sigma > 0$ is uniquely determined by the following implicit relation,
\[
\chi_N - \frac{\sigma}{1 - (\epsilon \sigma)^2} = \chi_S \frac{\sigma}{\sqrt{4 D_S \alpha + \sigma^2}}.
\]
\end{theorem}

For the proof of this theorem we assume that the chemoattractant $S$ has a unique maximum at $z = 0$ and that the nutrient $N$ is increasing. Then, this initial ansatz allows us to decouple the system of equations \eqref{basicmodel}, enabling us to solve for $\rho$, and consequently for $S$. Finally, we will validate the ansatz, which will lead to the equation for $\sigma$ as stated in Theorem \ref{thmacro}.

\paragraph{Kinetic travelling pulse} It is natural to expect that some of the qualitative behaviours observed at the macroscopic scale, such as the existence of travelling pulse solutions, should also be present at the mesoscopic level. This motivates the study the existence of travelling wave solutions directly within the mesoscopic framework, aiming to confirm that such wave-like structures are inherent to the dynamics of bacterial motion, regardless of the scale under consideration.
In the following, we summarise the results in \cite{calvez2016chemotactic} for completeness.

We are going to denote the wave speed as $c$, and not as $\sigma$ like in the previous sections. The space variable of the moving pulse at speed $c$ is denoted by $z = x-ct$.
We use the following convention: the superscripts $^+$ and $^-$ denote the position of the velocity $v$ with respect to the speed $c$: for $v>c$ and for $v<c$, respectively. Whereas the subscripts $_+$ and $_-$ denote the relative position with respect to the origin: for $z>0$ and for $z<0$, respectively. 


Similarly, we denote the spatial density of the same group of bacteria $f^-_+(z,v)$ as 
\[
    \rho^-_+(z) = \int_{v<c} f_+(z,v)d\nu(v).
\]
Moreover, the tumbling rate can take only four values $1\pm\chi_S\pm\chi_N$ denoted by $T^+_+, T^-_+, T^+_-, T^-_-$ which depend on weather $v$ crosses $c$ or weather $z$ crosses the origin or not. We illustrate this convention in Figure \ref{tumbling rates} in Appendix \ref{app: T}.




Therefore, $T\in[1-\chi_S-\chi_N, 1+\chi_S+\chi_N]$ and we conclude that the tumbling frequency is higher when bacteria move in the favourable direction and smaller when bacteria move in the unfavourable direction.

\begin{theorem}[\cite{calvez2016chemotactic}] \label{theorem_existence_mesos}
 Assume $(\chi_S, \chi_N) \in (0, 1/2) \times [0, 1/2)$. Under the assumption (A1), and conditions 
\begin{equation} 
    \frac{c_* + \sqrt{(c_*)^2 + 4 \alpha D_S}}{c_* + \sqrt{(c_*)^2 + 4 \alpha D_S} + 2D_S \lambda_-(c_*)} \le \text{explicit constant},
    \label{constraint1}
\end{equation}

\begin{equation}
    \frac{-c^* + \sqrt{(c^*)^2 + 4 \alpha D_S}}{-c^* + \sqrt{(c^*)^2 + 4 \alpha D_S} + 2D_S \lambda_+(c^*)} \le \text{explicit constant},
    \label{constraint2}
\end{equation}

\begin{equation}
    \text{Either } \chi_N \ge \chi_S \hspace{0.5cm} \text{or} \hspace{0.5cm} \frac{\sqrt{\alpha / D_S} + \lambda_+ (0)}{\sqrt{\alpha / D_S} + \lambda_- (0)} \le \text{explicit constant},
    \label{constraint3}
\end{equation}

where $\lambda_-(0)$ and $\lambda_+(0)$
 are the smallest positive roots of, respectively 
 \begin{equation}
     \int \frac{v}{T_-(v) + \lambda v} d\nu (v) = 0, \hspace{0.4cm} \int \frac{v}{T_+(v) - \lambda v} d \nu (v) = 0.
     \nonumber
 \end{equation}
 There exist a non negative c, and positive functions $(f, S, N) \in (L^1 \cap L^\infty (\mathbb{R} \times V)) \times \mathcal{C}^2(\mathbb{R}) \times \mathcal{C}^2 (\mathbb{R})$ solution of the travelling wave problem,
 \begin{equation}
         \left \{
    \begin{aligned}
       (v-c) \partial_z f(z,v) = \int T(z, v'-c)f(z,v') d\nu (v') - T(z, v-c)f(z,v),\\
       -c \partial_z S(z) - D_S \partial_z^2 S(z) + \alpha S(z) = \beta \rho (z),\\
       -c \partial_z N(z) - D_N \partial_z^2 N(z) = - \gamma \rho (z) N(z),
    \end{aligned}
  \right .
  \label{eq_sys_mesoscopic}
\end{equation}
where the tumbling rate $T$ is given by 

\[
    T(z, v-c) = 1 - \chi_S \sign((v-c)\partial_zS(z)) - \chi_N \sign((v-c)\partial_zN(z)) .
\]

\end{theorem}

The details of the proof of this theorem can be found in \cite{calvez2016chemotactic}.

\noindent An important issue discussed in \cite{calvez2018well} is regarding the uniqueness of the travelling wave, where the authors numerically observed that travelling wave exists for several values of $c$, in a parameter regime which violates conditions \eqref{constraint1}-\eqref{constraint3}. This clearly contradicts the results observed at the macroscopic level, where there exists a unique wave speed.  

\textcolor{red}{(this is something we should look into in our AP scheme)}







\section{Micro-macro decomposition}

To design an asymptotic preserving scheme which automatically preserves the macroscopic limit, a micro-macro formulation needs to be derived. 
We decompose the solution $f(t,x,v)$ as 
\begin{equation}\label{def:f_decompose}
    f(t,x,v) = \rho(t,x) + \varepsilon g(t,x,v)\ .
\end{equation}
Substituting the micro-macro decomposition of $f(t, x, v)$ given by \eqref{def:f_decompose} into the  kinetic model \eqref{eq: inital model with scaling}, integrating over $v$ and noticing the fact that $\langle T_\varepsilon \rangle=1 \, \textnormal{and}\, \langle f\rangle = \rho$, we have the equation for the macroscopic quantity $\rho(t, x)$
\begin{equation*}
    \partial_t \rho +  \langle v \nabla  g\rangle  =  0\ .
\end{equation*}
To get the equation for $g$, we use the projection technique. 
For simplicity of notations, we introduce the projection operator $\Pi$ defined as 
\begin{equation*}
    \Pi f(t,{v},{x}) = \langle f(t,{v},{x})\rangle \ .
\end{equation*}
It is easy to check that, {for $I$ the identity operator,} 
\begin{align*}
    & (I - \Pi) f = \varepsilon(I - \Pi) g = \varepsilon g\ ,\\
    & (I - \Pi) \left(\int_VT_\varepsilon(v')f^\varepsilon(v')\diff v'-T_\varepsilon(v)f^\varepsilon(v)\right) = -\varepsilon\left(g+\rho \psi_1 \right) \ .
\end{align*}
Finally taking the operator $I-\Pi$ into equation \eqref{eq: inital model with scaling}, we get
\begin{align*}
    &\partial_t g + \frac{1}{\varepsilon} (I-\Pi)\left[v \partial_x g \right] 
    = \frac{1}{\varepsilon^2}\left[-(v\partial_x\rho-\rho \psi_1  - g\right] \ .
\end{align*}

As a summary, by decomposing $f$ as \eqref{def:f_decompose}, the following micro-macro formulation of the system is derived
\begin{align}\label{eq:macro_micro}
\begin{cases}
	&\partial_t \rho +  \langle v\partial_x g\rangle  = 0\ , \\
	&\partial_t g + \frac{1}{\varepsilon} (I-\Pi) ({v}\partial_x g) = \frac{1}{\varepsilon^2}(-({v}\partial_x \rho-\rho\psi_1 -  g) \ ,\\
	&  \partial_t S = \tilde{D}_S\partial_{xx} S -\tilde{\alpha}S+\tilde{\beta}\rho\ ,\\
 & \partial_tN = \tilde{D}_N\partial_{xx} N-\tilde{\gamma}\rho N\ .
\end{cases}
\end{align}

\section{An asymptotic preserving finite difference scheme}

Here we introduce an asymptotic preserving scheme by discretising \eqref{eq:macro_micro} using a finite difference method. 
We consider the one-dimensional case for simplicity, i.e $x,v\in[x_{\textnormal{min}},x_{\textnormal{max}}]\times[v_{\textnormal{min}},v_{\textnormal{max}}]$ with homogeneous Neumann boundary conditions in both variables.
We denote  $x_{j+\frac{1}{2}}=(j+\frac{1}{2})\Delta x$ and $x_{j}=j\Delta x$ for $j=0,..,N_x-1$,  $v_k=k\Delta v$ for $ k=1,...,N_v$ and $t_n=n\Delta t$ for $n\geq 0$, where $ N_x=(x_{\textnormal{max}}-x_{\textnormal{min}})/\Delta x$ and $N_v=(v_{\textnormal{max}}-v_{\textnormal{min}})/\Delta v$. We consider the approximation of the unknown density function  $\rho(t,x)$ computed on half grid points $x_{j+\frac{1}{2}}$ with $\rho_{j+\frac{1}{2}}^n\approx\rho(t_n,x_{j+\frac{1}{2}})$, while the unknown perturbation function $g(t,x,v)$ is computed on  $x_{j}$ with $g_{j,k}^n\approx g(t_n,x_{j},v_k)$. %Also,   $g(t_n,x_{j+\frac{1}{2}},v_k)$ is calculated by interpolation:
%$$g(t_n,x_{j+\frac{1}{2}},v_k)\approx \bar{g}_{j+\frac{1}{2},k}^n:=(g_{j,k}^n+g_{j+1,k}^n)/2\ .$$

For simplicity of notations, we further introduce the standard finite difference operators $\delta_t^+$ and $\delta_x$, which are numerical approximations of $\partial_t$ and $\partial_x$, respectively, and defined as 

$$\delta_t^+\rho_{j+\frac{1}{2}}^n=\frac{\rho_{j+\frac{1}{2}}^{n+1}-\rho_{j+\frac{1}{2}}^n}{\Delta t}, \quad\delta_xg_{j+\frac{1}{2},k}^n=\frac{g_{j+1,k}^n-g_{j,k}^n}{\Delta x},\quad \delta_x\rho_{j}^n=\frac{\rho_{j+\frac12}^n-\rho_{j-\frac12}^n}{\Delta x}\ .$$

The composite of two operators $\delta_x$, which is denoted as $\delta_x^2$, is then defined to be 
$$
\delta_x^2\rho_{j+\frac{1}{2}}^n=\frac{\delta_ x\rho_{j+1}^n-\delta_x\rho_{j}^n}{\Delta x}=\frac{\rho_{j+\frac{3}{2}}^n-2\rho_{j+\frac{1}{2}}^n+\rho_{j-\frac{1}{2}}^n}{(\Delta x)^2}
$$
which is the numerical approximation of $\partial_x^2$. Another important notation to be introduced is $\langle \cdot \rangle_h$, which is defined as
$$
\langle \eta_{j,k}^n \rangle_h:=\Delta v\sum_k\eta_{j,k}^n\ ,
$$
where $\eta_{j,k}^n\approx \eta(t_n,x_j,v_k)$ for some general function $\eta(t,x,v)$. Obviously, $\langle \eta_{j,k}^n \rangle_h$ is the finite difference approximation of $\langle \eta(t_n,x_j,v))\rangle:=\int_V \eta(t_n,x_j,v) dv$


\iffalse
 Finally, to better approximate $\rho$ at $x=x_{j+\frac{1}{2}}$ and $t=t_n$, we introduce the notation $\Phi_{j+\frac{1}{2}}^{n_1,n}$, which is defined as 
 \begin{eqnarray}
     \Phi_{j+\frac{1}{2}}^{n_1,n}=\left\{ \begin{array}{rc}\rho_j^{n_1} ,& \mbox{if} \quad  \delta_xS_{j+\frac{1}{2}}^n \geq 0,\\
     \rho_{j+1}^{n_1} ,  & \mbox{if} \quad \delta_xS_{j+\frac{1}{2}}^n < 0,\end{array}\right.
 \end{eqnarray}
 where $n_1=n$ or $n+1$. As shown in [] $\Phi_{j+\frac{1}{2}}^{n_1,n}$ approximates $\rho$ at $t=t_n$ and $x=x_{j+\frac{1}{2}}$ in an upwind manner and thus help improve the stability of the numerical scheme.
 \fi
 
  With the notations defined, the system \eqref{eq:macro_micro} can be discretised as
   \begin{equation}
\begin{cases}\label{eq: main scheme}
   &\delta_t^+\rho_{j+\frac{1}{2}}^n +\langle v_k \delta_x (g_{j+\frac12,k}^{n+1})\rangle_h=0\ ,\\
   
        &\delta_t^+ g_{j,k}^n + \frac{1}{\epsilon}(I-II_h)M_{j,k}^n =\frac{1}{\epsilon^2}(K_{j,k}^{n,n+1}-g_{j,k}^{n+1})\ ,\\
        
        &\delta_t^+S_{j+\frac{1}{2}}^{n+1} = \tilde{D}_S \delta _x^2S_{j+\frac{1}{2}}^{n}-\tilde{\alpha} S_{j+\frac{1}{2}}^{n}+\tilde{\beta}\rho_{j+\frac{1}{2}}^{n}\ ,\\
        
     &\delta_t^+N_{j+\frac{1}{2}}^{n+1} = \tilde{D}_N \delta _x^2N_{j+\frac{1}{2}}^{n}-\tilde{\gamma} \rho_{j+\frac{1}{2}}^{n}N_{j+\frac{1}{2}}^{n}\ ,\\
   \end{cases}
   \end{equation}
where
\begin{align}
M_{j,k}^n&=(v_k^+\delta_x(g_{j-\frac{1}{2},k}^n)-v_k^-\delta_x(g_{j+\frac{1}{2},k}^n))\ ,\\
K_{j,k}^{n,n+1}
& = -v_k\delta_x\rho_{j}^{n+1}+\rho^{[S],n+1}_{j}\phi_1[S_{j}^n](v_k)+\rho^{[N],n+1}_{j}\phi_2[N_{j}^n](v_k)
\end{align}
for $v^+=\max\{v,0\}$ and $v^-=\max\{-v,0\}$ and 
{
\[
\rho^{[S],n+1}_{j}=\begin{cases}
     \rho^{n+1}_{j-\frac{1}{2}}\ ,  &\textnormal{if}\ \  \delta_x S_{j}^n\geq 0\ ,\\
    \rho^{n+1}_{j+\frac{1}{2}}\ , &\textnormal{if}\ \ \delta_x S_{j}^n< 0\ ,
\end{cases}
\quad \rho^{[N],n+1}_{j}=\begin{cases}
     \rho^{n+1}_{j-\frac{1}{2}}\ , & \textnormal{if}\ \  \delta_x N_{j}^n\geq 0\ ,\\
    \rho^{n+1}_{j+\frac{1}{2}}\ , & \textnormal{if}\ \ \delta_x N_{j}^n< 0\ .
    \end{cases}
\]
Following the idea in \cite{lemou2008new}, the scheme \eqref{eq: main scheme} can be solved efficiently. Instead of solving the system \eqref{eq: main scheme} directly, where all densities $\rho_{j+\frac{1}{2}}^{n+1}$ and perturbations $g_{j,k}^{n+1}$ are coupled so that a large linear system needs to be inverted, we introduce $\tilde{g}_{j,k}^{n+1}$, which satisfies
\begin{eqnarray}\label{eqg}
    \frac{\tilde{g}_{j,k}^{n+1}-g_{j,k}^n}{\Delta t}+\frac{1}{\epsilon}(I-II)M_{j,k}^n=\frac{1}{\epsilon^2}(\widetilde{K}_{j,k}^{n}-\tilde{g}_{j,k}^{n+1})\ ,
\end{eqnarray}
where
$$
\widetilde{K}_{j,k}^{n}=-v_k\delta_x\rho_{j}^{n}+\rho^{[S],n}_{j}\phi_1[S_{j}^n](v_k)+\rho^{[N],n}_{j}\phi_2[N_{j}^n](v_k)\ .
$$
By reformulating \eqref{eqg}, it is easy to see that
\iffalse
\begin{eqnarray}\label{gtilda}
\left( \frac{1}{\Delta t}+\frac{1}{\epsilon ^2}\right)\tilde{g}_{j+\frac{1}{2},k}^{n+1}=\frac{g_{j+\frac{1}{2},k}^n}{\Delta t}-\frac{1}{\epsilon}(I-II)\left(v_k^+\delta_x(g_{j,k}^n)-v_k^-\delta_x(g_{j+1,k}^n)\right)\nonumber\\
+\frac{1}{\epsilon^2}\left( -v_k\psi_0(v_k)\delta_x\rho_{j+\frac{1}{2}}^{n+1}+\psi_1(v_k,\delta_x S_{j+\frac{1}{2}}^n)\Phi_{j+\frac{1}{2}}^{n+1,n}\right)
\end{eqnarray}
\fi
 all the unknowns $\tilde{g}_{j,k}^{n+1}$ can be solved explicitly. By comparing \eqref{eqg} and the second equation in \eqref{eq: main scheme}, it can be observed that 
\begin{align}\label{g}
    g_{j,k}^{n+1}
    &=\tilde{g}_{j,k}^{n+1}+\frac{1}{\epsilon^2}\left(\frac{1}{\Delta t}+\frac{1}{\epsilon^2}\right)^{-1}\left( K_{j,k}^{n,n+1}-\widetilde{K}_{j,k}^n \right)\ .
\end{align}
Then by substituting \eqref{g} into first equation \eqref{eq: main scheme}, a system which contains only the unknowns for the densities is derived. Specifically, we have  
\begin{eqnarray}\label{rho}
  \frac{\rho_j^{n+1}}{\Delta t}-\delta_x(a_*^n\delta_x\rho_*^{n+1})_j  +\delta_x(b_*^n\rho^{[S],n+1}_{j})_j 
  +\delta_x(c_*^n\rho^{[N],n+1}_{j})_j
  =r_j 
\end{eqnarray}
where the coefficients $a_{j+\frac{1}{2}}^n, b_{j+\frac{1}{2}}^n, c_{j+\frac{1}{2}}^n$ and residuals $r_j^n$ can be explicitly computed via
\begin{eqnarray*}
    a_{j+\frac{1}{2}}^n&=&\frac{\Delta t}{\epsilon^2+\Delta t} D_\rho\ ,\nonumber\\
    b_{j+\frac{1}{2}}^n&=&\frac{\Delta t}{\epsilon^2+\Delta t} \langle v_k\phi_1[S_{j}^n](v_k)\rangle_h\ ,\nonumber\\
    c_{j+\frac{1}{2}}^n&=&\frac{\Delta t}{\epsilon^2+\Delta t} \langle v_k\phi_2[N_{j}^n](v_k)\rangle_h\ ,\nonumber\\
    r_j^n&=&\frac{{\rho_j^{n}}}{\Delta t}-\langle v_k\delta_x(\tilde{g}_{*,k}^{n+1})_j\rangle_h
    -\delta_x(a_*^n\delta_x\rho_*^n)_j
    +\delta_x(b_*^n\rho^{[S],n}_{j})
    +\delta_x(c_*^n\rho^{[N],n}_{j})\ ,
    \end{eqnarray*}
    where $D_\rho\coloneqq\langle v_k^2\rangle_h$ is a diffusion coefficient, analogously defined in \eqref{eq: moment}. 
To solve all the unknowns $\rho_j^{n+1}$ from the system \eqref{rho}, only a tridiagonal matrix needs to be inverted. And the unknowns 
 $g_{j+\frac{1}{2},k}^{n+1}$ can be solved explicitly via \eqref{g}. In this way, we efficiently update the system \eqref{rho} from $t=t_n$ to $t=t_{n+1}$}


{
The numerical approximations $\rho_{j+\frac12}^{n+1}$ and $g_{j,k}^{n+1}$ obtained from \eqref{eq: main scheme} are, in general, not automatically positivity preserving.
In practice, however, small negative undershoots may appear and may subsequently lead to numerical instabilities.
To ensure that the distribution
$$
f_{j,k}^{n+1} = \rho_j^{n+1} \psi_0(v_k) + \varepsilon g_{j,k}^{n+1}
$$
remains nonnegative and physically meaningful, we apply a {\GER Zhang–Shu type flux limiter \cite{ZS-limiter}}.
The limiter adjusts only the microscopic fluctuation to enforce positivity, while strictly preserving the total mass. 
A detailed description of the limiting procedure is provided in Appendix~\ref{sec:limiter}.
}

\subsection{Asymptotic preserving property}
In this section we formally show the asymptotic preserving property of the numerical scheme, where we let $\epsilon\to 0$ in \eqref{eq: main scheme}. We aim to show that the numerical scheme for the kinetic equation \eqref{eq: inital model with scaling} converges, as $\epsilon\to 0$, to the numerical scheme of the macroscopic equation \eqref{basicmodel} (first equation).

{Formally, by taking $\epsilon\to0$ in the second equation of \eqref{eq: main scheme}, we have 
\begin{align*}
    g_{j,k}^{n+1} & = K_{j,k}^n \\
    & = -v_k\delta_x\rho_{j}^{n+1}+
    \rho^{[S],n+1}_{j}\phi_1(v_k) \sign(\delta_x S_j^n) +\rho^{[N],n+1}_{j} \phi_2(v_k) \sign(\delta_x N_j^n)\ .
\end{align*}
A direct computation shows that
\begin{align}
    \langle v_k \delta_x(g_{j+\frac12,k}^{n+1}) \rangle_h = &
    -\langle v_k^2  \rangle_h \delta^2_x \rho_{j+\frac12}^{n+1}\nonumber \\
    &+ \langle v_k \phi_1(v_k) \rangle_h \frac{\sign(\delta_x S_{j+1}^n) \rho^{[S],n+1}_{j+1} - \sign(\delta_x S_{j}^n) \rho^{[S],n+1}_{j}}{\Delta x}\nonumber\\
    &+\langle v_k \phi_2(v_k) \rangle_h \frac{\sign(\delta_x N_{j+1}^n) \rho^{[N],n+1}_{j+1} - \sign(\delta_x N_{j}^n) \rho^{[N],n+1}_{j}}{\Delta x}\ .\label{eq: AP computation}
\end{align}
}
{
Substituting the above result \eqref{eq: AP computation} into the first equation in  \eqref{eq: main scheme} and introducing $D_\rho=\langle v_k^2\rangle_h$, 
$\chi_S \coloneqq \langle v_k \phi_1(v_k) \rangle_h$ and 
$\chi_N \coloneqq \langle v_k \phi_2(v_k) \rangle_h$, 
we finally get the finite difference scheme 
\begin{align*}
    \delta_t^+\rho_{j+\frac{1}{2}}^n = &D_\rho \delta_x^2 \rho_{j+\frac12}^{n+1} \\
    &- \chi_S \frac{\sign(\delta_x S_{j+1}^n) \rho^{[S],n+1}_{j+1} - \sign(\delta_x S_{j}^n) \rho^{[S],n+1}_{j}}{\Delta x} \\
    &- \chi_N \frac{\sign(\delta_x N_{j+1}^n) \rho^{[N],n+1}_{j+1} - \sign(\delta_x N_{j}^n) \rho^{[N],n+1}_{j}}{\Delta x},
\end{align*}
which is indeed the scheme for solving the macroscopic model. 
In this way, we verified the asymptotic preserving property of our scheme \eqref{eq: main scheme}.
}
\subsection{Bound-preserving property}
{In general, we do not have the bound-preserving property for the kinetic model {\GER(shall we say why not, basically because we cannot prove that $r_j^n\geq 0$)} {\GER (comment that the flux limiter imposes a lower bound on the solution)}. However, the property can be proved rigorously for the limiting scheme as $\epsilon\to 0$. 
For simplicity, we denote $c_j^n = \sign(\delta_x S_{j}^n)$, $d_j^n = \sign(\delta_x N_{j}^n)$, then the scheme can be rewritten as 
\begin{align*}
    \delta_t^+\rho_{j+\frac{1}{2}}^n &= D_\rho \delta_x^2 \rho_{j+\frac12}^{n+1} 
    - \chi_S \frac{c_{j+1}^n \rho^{[S],n+1}_{j+1} - c_{j}^n \rho^{[S],n+1}_{j}}{\Delta x} 
    - \chi_N \frac{d_{j+1}^n \rho^{[N],n+1}_{j+1} - d_{j}^n \rho^{[N],n+1}_{j}}{\Delta x}\ .
\end{align*}

\begin{proposition}[Mass conservation, positivity, and upper bound]
Consider the finite-difference scheme for the macroscopic density 
$\rho_j^n$,
\[
M^n \brho^{n+1} = \brho^n ,
\]
where $M^n$ is the tridiagonal matrix with entries
\begin{align*}
m_{j,j-1}^n &= - D_\rho \frac{\Delta t}{(\Delta x)^2}
             - \chi_S \frac{\Delta t}{\Delta x} (c_j^n)^{+}
             - \chi_N \frac{\Delta t}{\Delta x} (d_j^n)^{+}, \\
m_{j,j}^n   &= 1 + 2 D_\rho \frac{\Delta t}{(\Delta x)^2}
             + \chi_S \frac{\Delta t}{\Delta x} \left[(c_{j+1}^n)^{+} + (c_j^n)^{-}\right]
             + \chi_N \frac{\Delta t}{\Delta x} \left[(d_{j+1}^n)^{+} + (d_j^n)^{-}\right], \\
m_{j,j+1}^n &= - D_\rho \frac{\Delta t}{(\Delta x)^2}
             - \chi_S \frac{\Delta t}{\Delta x} (c_{j+1}^n)^{-}
             - \chi_N \frac{\Delta t}{\Delta x} (d_{j+1}^n)^{-},
\end{align*}
where 
\[c_j^n = \operatorname{sign}(\delta_x S_j^n),
\qquad
d_j^n = \operatorname{sign}(\delta_x N_j^n),
\]
and $(\cdot)^{+}$ and $(\cdot)^{-}$ denote the positive and negative parts, respectively.

Assume:
\begin{enumerate}
    \item $D_\rho > 0$, and $\chi_S, \chi_N \ge 0$.
    \item The CFL-type condition
    \begin{equation}\label{eq:CFL}
    \frac{\Delta t}{\Delta x} \le \frac{1}{4(\chi_S + \chi_N)}
    \end{equation}
    holds.
    \item The initial data satisfy $\rho_j^0 \ge 0$ for all $j$.
\end{enumerate}
Then the numerical solution satisfies, for all $n \ge 0$:
\begin{enumerate}
    \item {Mass conservation:} ${\|\brho^{n+1}\|_1 = \|\brho^{n}\|_1}$.
    \item {Positivity preservation:} If $\rho_j^n \ge 0$ for all $j$, then 
    $\rho_j^{n+1} \ge 0$ for all $j$.
    \item {Uniform upper bound:} Let $\|\brho^0\|_1 = \sum_j \rho_j^0$ and 
    $\bar\rho := 2\|\brho^0\|_1$. Then
    \[
    0 \le \rho_j^n \le \bar\rho, \qquad \forall j, \ \forall n.
    \]
\end{enumerate}
Hence, the scheme is mass conserving, positivity preserving, and uniformly bounded in $L^\infty$ under the stated CFL condition.
\end{proposition}





















% The numerical scheme can be reformulated in matrix form 
% $$
%     M^n \rho^{n+1} = \rho^{n},
% $$
% where $M^n$ is a tridiagonal matrix with
% \begin{align*}
%     & m_{j, j-1} = -D_\rho \frac{\Delta t}{(\Delta x)^2} 
%     - \chi_s \frac{\Delta t}{\Delta x} (c_{j}^n)_+ 
%     - \chi_N \frac{\Delta t}{\Delta x} (d_{j}^n)_+, \\
%     & m_{j,j} = 1 + 2D_\rho \frac{\Delta t}{(\Delta x)^2} 
%     + \chi_s \frac{\Delta t}{\Delta x} \left[(c_{j+1}^n)_+ + (c_j^n)_-\right] 
%     + \chi_N \frac{\Delta t}{\Delta x} \left[(d_{j+1}^n)_+ + (d_j^n)_-\right], \\
%     & m_{j, j+1} = -D_\rho \frac{\Delta t}{(\Delta x)^2} 
%     - \chi_s \frac{\Delta t}{\Delta x} (c_{j+1}^n)_- 
%     - \chi_N \frac{\Delta t}{\Delta x} (d_{j+1}^n)_-. \\
% \end{align*}
% Here we use $(c)_+$ and $(c)_-$ to denote the positive part and negative part of $c$, respectively. Obviously, a direct computation shows that 
% \begin{equation*}
%     m_{j-1,j} + m_{j,j} + m_{j+1,j} = 1.
% \end{equation*}
% Therefore, the following properties hold.

% \begin{enumerate}
%     \item (Mass conservation) $\mathbf{1}^T \rho^{n+1} = \mathbf{1}^T M^n \rho^{n+1} = \mathbf{1}^T \rho^{n}$.
%     \item (Positive-preserving) If $\rho^n \ge 0$, then we have $\rho^{n+1} \ge 0$.
    
% It immediately leads to 
% $$
%     \|\rho^0\|_1 = \|\rho^n\|_1, \, n=1,2,\dots.
% $$
% \item (Upper bound) As long as 
% \begin{equation}\label{eq:CFL}
%     \frac{\Delta t}{\Delta x} \le \frac{1}{4(\chi_S + \chi_N)},  
% \end{equation}
% we have $\rho_j^n \le 2\|\rho\|_1$ for all $j$.
% \end{enumerate}
% Generally speaking, if we wish to have the bound-preserving property, i.e. when $\rho^n \le \bar{\rho}$, we have $\rho^{n+1}\le \bar{\rho}$, we need more conditions. 
% One possible way is to choose an upper bound $\bar{\rho}$ such that 
% $$
%     (m_{j,j-1} + m_{j,j} + m_{j,j+1}) \bar{\rho} \ge \rho_j^{n}, 
% $$
% holds true for all $j = 1,2,\dots,N-1$. 
% Or equivalently 
% $$
%     \left[1 + \chi_S \frac{\Delta t}{\Delta x}((a_{j+1})_+ -(a_{j+1})_- + (a_j)_- - (a_j)_+) + \chi_N \frac{\Delta t}{\Delta x}((b_{j+1})_+ -(b_{j+1})_- + (b_j)_- - (b_j)_+)\right]\bar{\rho} \ge \rho_j^n.
% $$
% Or 
% $$
%     \left[1 + \chi_S \frac{\Delta t}{\Delta x}(a_{j+1} - a_j) + \chi_N \frac{\Delta t}{\Delta x}(b_{j+1} - b_j)\right]\bar{\rho} \ge \rho_j^n.
% $$
% Noticing that $a_{j+1}-a_j \in (-2, 2)$ and $b_{j+1}-b_j \in (-2, 2)$, when
% $$
%     \frac{\Delta t}{\Delta x} \le \frac{\bar{\rho} - \rho_j^n}{2(\chi_S + \chi_N)\bar{\rho}},
% $$
% we have 
% $\rho^{n+1}\le \bar{\rho}$. (Remark: The time step constraint is time dependent and we need $\rho_j^n$ is not close to $\bar{\rho}$. )
% For example, we can take $\bar{\rho} = 2\|\rho^0\|_1$, then \begin{align*}
%     \bar{\rho} - \rho_j^n & =  2\|\rho^0\|_1 - \rho_j^n \\
%     & = 2\|\rho^n\|_1 - \rho_j^n \\
%     & \ge \|\rho^n\|_1 = \|\rho^0\|_1.
% \end{align*}
% Here we used the fact that $\|\rho\|_\infty \le \|\rho\|_1$.
% So a uniform constraint would be 
% \begin{equation*}
%     \frac{\Delta t}{\Delta x} \le \frac{1}{4(\chi_S + \chi_N)}.
% \end{equation*}

\begin{proof}
We begin by observing that the tridiagonal matrix $M^n$ satisfies
\[
m_{j,j-1}^n + m_{j,j}^n + m_{j,j+1}^n
= 1 + \chi_S \frac{\Delta t}{\Delta x} (c_{j+1}^n - c_j^n)
      + \chi_N \frac{\Delta t}{\Delta x} (d_{j+1}^n - d_j^n)
=: F_j^n .
\]
From the bounds on $c_{j+1}^n - c_j^n$ and $d_{j+1}^n - d_j^n$, since  $c_j^n, d_j^n \in \{-1,0,1\}$, we have 
\[
F_j^n \ge 1 - 2(\chi_S+\chi_N)\frac{\Delta t}{\Delta x}.
\]
Under the CFL condition \eqref{eq:CFL} we obtain the uniform lower bound $F_j^n \ge \frac{1}{2}$.


To prove the three properties, we proceed in steps.\\

\noindent\underline{1. Mass conservation.}
{Noticing the fact that $\mathbf{1}^T M^n = \mathbf{1}^T$, a direct computation shows that  
$$
\|\brho^{n+1}\|_1 = \mathbf{1}^T \brho^{n+1} = \mathbf{1}^T M^n  \brho^n = \mathbf{1}^T \brho^{n} = \|\brho^n\|_1.
$$}

\vspace{0.2cm}
\noindent\underline{2. Positivity preservation.}
The matrix $M^n$ is an M-matrix since all off-diagonal entries $m_{j,j-1}^n$ and $m_{j,j+1}^n$ are nonpositive, the diagonal entries satisfy $m_{j,j}^n \ge F_j^n \ge \frac{1}{2}$ and $M^n$ is strictly diagonally dominant by rows, i.e. $m_{j,j}^n \ge |m_{j,j-1}^n| + |m_{j,j+1}^n| + \frac{1}{2}$.
Thus $M^n$ is a nonsingular M-matrix, and it is well known that its inverse is nonnegative
Therefore,
\[
\brho^{\,n+1} = (M^n)^{-1} \brho^n \ge 0
\]
whenever $\brho^n \ge 0$. This proves positivity preservation.

\vspace{0.2cm}
\noindent\underline{3. Uniform upper bound.}
Let  $\bar\rho := 2\|\brho^0\|_1$.
Since the scheme preserves mass, we have for all $n$, that $\|\brho^n\|_1 = \|\brho^0\|_1$, and therefore $\rho_j^n \leq \|\brho^n \|_\infty\leq \|\brho^n\|_1 = \|\brho^0\|_1 = \frac12 \bar\rho$.


Next, using $F_j^n \ge 1/2$, we compute
\[
(M^n (\bar\rho \mathbf{1}))_j
= \sum_{k=j-1}^{j+1} m_{j,k}^n \bar\rho
= F_j^n \bar\rho
\ge \frac12 \bar\rho
\ge \rho_j^n.
\]
Hence
\[
M^n (\bar\rho \mathbf{1}) \ge \brho^n.
\]
Applying the inverse (which is nonnegative),
\[
\bar\rho \mathbf{1} = (M^n)^{-1} M^n(\bar\rho \mathbf{1})
\ge (M^n)^{-1} \brho^n = \brho^{n+1}.
\]
Thus $0 \le \rho_j^{n+1} \le \bar\rho$ for all $j$.
Combining the three steps yields the stated result.
\end{proof}








% Proof: For simplicity, we denote 
% \begin{equation*}
%     F_j := m_{j,j-1} + m_{j,j} + m_{j,j+1}
%         = 1 + \chi_S \frac{\Delta t}{\Delta x}(c_{j+1} - c_j)
%           + \chi_N \frac{\Delta t}{\Delta x}(d_{j+1} - d_j).
% \end{equation*}
% From the bounds on $c_{j+1}-c_j$ and $d_{j+1}-d_j$, we have
% \[
%     F_j \ge 1 - 2(\chi_S + \chi_N)\frac{\Delta t}{\Delta x}.
% \]
% Under the CFL condition~\eqref{eq:CFL}, it follows that $F_j \ge \tfrac{1}{2}$.
% Let $\bar{\rho} := 2\|\rho^0\|_1$. 
% By the discrete inequality $\|\rho\|_\infty \le \|\rho\|_1$ and the mass conservation, we obtain
% \[
%     \rho_j^n \le \|\rho^n\|_\infty \le \|\rho^n\|_1 = \|\rho^0\|_1.
% \]
% Consequently, for each $j$,
% \[
%     F_j \bar{\rho} 
%     \ge \tfrac{1}{2} \cdot 2\|\rho^0\|_1
%     = \|\rho^0\|_1
%     \ge \rho_j^n.
% \]
% Hence $M(\bar{\rho}\, \mathbf{1}) = (F_j \bar{\rho})_j \ge \rho^n$. 
%  By the monotonicity of $M^{-1}$ (since $M$ is an $M$-matrix), we have
% \[
%     \rho^{n+1} = M^{-1}\rho^n \le M^{-1}M(\bar{\rho}\, \mathbf{1}) = \bar{\rho}\, \mathbf{1}.
% \]
% This completes the proof.
% }

\section{Numerical experiments}

% \textcolor{blue}{Compute the travelling wave solution speed at the macro level for different values of $\chi_N$ and $\chi_S$, see reference [8]. We can compare to the numerical velocity obtained at the kinetic level for the same parameters.}

% \textcolor{red}{(To Xinran: Please explain here the details, the choices of the functions, the parameter values and include some plots for the macroscopic and kinetic travelling waves.)}

% \textcolor{red}{
% \[(\chi_N-\sigma)^2(4D_S\alpha+\sigma^2)=\chi_S^2\sigma^2\]
% \[ \chi_N=1.1,\quad D_S=2,\quad \alpha=0.4,\quad \chi=\frac{1}{2}\]
% \[\sigma^4-2.2\sigma^3+4.16\sigma^2-7.04\sigma+3.872=0\]
% \[\sigma=0.8794\]
% From the data observed in the previous experiments in [refs] we can compute the following
% \[
% A=1\ , \ \  B=1/2\ ,\ \ C=1.1\ ,\ \ D=2\ ,\ \ E=0.4\ , \ \ F=4\cdot 10^7\textnormal{cells}^{-1}\ ,\ \ G=2\ ,\ \ H=4\cdot 10^7\textnormal{cells}^{-1}\ .
% \]}

\textcolor{red}
{\begin{itemize}
    \item Are we defining $\rho^\varepsilon_{\textnormal{kinetic}}$ and $\rho_{\textnormal{macro}}$?
    \item Define initial conditions for the kinetic solution.
    \item Specify space-time discretisation.
    % \item Check figure 3. Remove this figure
    % \item Make Table 1 a larger parameter study in a plot. We can do the same for the kinetic solution for a fixed epsilon. Vary $\chi_N$ and $\chi_S$ for numerical wave speed and analytical wave speed.
    % \item I don't see the difference between Table 2 and Table 3. Remove table 2.
\end{itemize}
}

In this section we present several numerical examples. In particular, we look for travelling wave solutions at both, the kinetic and macroscopic level, and compare both solutions when $\varepsilon\to 0$.

We consider a one-dimensional sufficiently large computational domain $[0,150]$ and employ a uniform spatial mesh with $\Delta x = 0.1$.
The homogeneous Neumann boundary condition is imposed at the left end of the domain, while a homogeneous Dirichlet boundary condition is applied at the right end.
For numerical simulations of the macroscopic model~\eqref{basicmodel}, 
we assume that cells initially concentrate on the left side of the domain,
with no chemoattractant signal and a constant nutrient level.  
the initial data are prescribed as
\[
\rho(0,x) = \rho_0 e^{-\lambda x},\qquad 
S(0,x) = 0,\qquad 
N(0,x) = N_0,
\]
for given constants $\rho_0,\ \lambda,\ N_0$.  
The remaining parameters in~\eqref{basicmodel} follow~\cite{saragosti2010mathematical}:
\[
D_\rho = \frac23,\qquad  
D_S = 2,\qquad 
D_N = 0,\qquad 
\alpha = 0.05,\qquad 
\beta = \gamma = 1.
\]
Figure \ref{fig:macro_analytic_cmp} shows the profiles of $\rho(t,x)$ with $\chi_S = 0.5$ and $\chi_N = 1.1$. 
In the figure, we showed that a travelling wave appears. 
Further, we showed that the shape of the profile is consistent with our analytical results \eqref{}. 
The travelling speed can be computed. 
By varying the values of $\chi_S$ and $\chi_N$, we should expect that the travelling speed $\sigma$ varies as well according to
$$
    \chi_N - \sigma = \chi_S \frac{\sigma}{\sqrt{4D_S\alpha + \sigma^2}}\ ,
$$
which by monotonicity arguments admits at most one solution. 
In Figure \ref{fig:travelling_speed}, we compared the numerical travelling speed with the analytical results for different values of $\chi_N$. 
When $\chi_S > \chi_N$, there is no travelling wave solution. When $\chi_S < \chi_N$, as shown in the figure, the numerical travelling speed is consistent with the analytical one. 



% In the absence of nutrient, there should be a stationary solution $\rho(x) = e^{-\frac{\chi_S}{D_\rho} x}$ as suggested in \cite{saragosti2010mathematical}. In other words, we should have $\sigma=0$, which is consistent with our result (\textcolor{red}{related to some figure?})
\begin{figure}
    \centering
    \includegraphics[width=0.5\linewidth]{profiles.eps}\includegraphics[width=0.5\linewidth]{compare_profile.eps}
    \caption{The figure on the left shows density profiles $\rho(t,x)$ at $T=20,60,100$. 
    The figure on the right compares the numerical density profile $\rho$ at $T=100$ with the analytical one \textcolor{red}{given as $\rho_\text{analytic} = \rho_0 e^{\lambda_n(x_p - x)_+ + \lambda_p(x-x_p)_+}$, where $\rho_0 = \frac{\Delta x \sum \rho_j}{1/\lambda_n + 1 / |\lambda_p|}$, with $\lambda_n = \frac{-\sigma +\chi_S + \chi_N}{D_\rho}$, $\lambda_p = \frac{-\sigma -\chi_S + \chi_N}{D_\rho}$, $\sigma$ being the travelling speed and $x_p$ being the position of the peak }.}
    \label{fig:macro_analytic_cmp}
\end{figure}
}




% \begin{table}[]
%     \centering
%     \begin{tabular}{c|c|c|c}
%         \hline
%         $\chi_N$ & 1.1 & 2 & 3 \\
%         \hline
%         Numerical travelling speed & 0.7230 & 1.5428 & 2.5210 \\
%         \hline
%         Analytical travelling speed & 0.7235 & 1.5376 & 2.5151 \\
%         \hline
%     \end{tabular}
%     \caption{Comparison of the numerical travelling speed and the analytical travelling speed. In the test, we fix $\chi_S = 0.5, D_s = 2, \alpha = 0.05$.
%     (For $\chi_N=3$, $\Delta x=0.01$,
%     $\chi_N=2$, $\Delta x=0.01$,
%     $\chi_N=3$, $\Delta x=0.02$)}
%     \label{tab:compare_speed_chiN}
% \end{table}


{ {
For the kinetic model \eqref{eq: inital model with scaling} \XR{(Is the reference of equation proper?)}, 
we adopt a constant equilibrium distribution 
$\psi_0(v) \equiv 1$ as in \cite{}, and choose the turning kernels
\[
\psi_1^{(S)}(v) = \frac{\chi_c}{\sqrt{2\pi}\sigma^3}\, v\, e^{-v^2/(2\sigma^2)}, 
\qquad 
\psi_1^{(N)}(v) = \frac{\chi_n}{\sqrt{2\pi}\sigma^3}\, v\, e^{-v^2/(2\sigma^2)},
\]
where $\sigma = 0.1$, and the parameters $\chi_c$ and $\chi_n$ are chosen to be the same as in the macro model.
% $$ 
%     \phi_1[S](v) = \phi_1[N](v) = \frac{v}{2\pi} e^{-\frac{v^2}{2}}
% $$
% Figure \ref{fig:conv_T} shows the profiles of \textcolor{red}{$\rho(t,x)$ (or $f(t,x,v)$ since the equation for $\rho$ is independent of $\varepsilon$)} with different values of $\varepsilon$.
}}
% \begin{figure}
%     \centering
%     \includegraphics[width=0.5\linewidth]{profiles_epsilon.eps}
%     \caption{We compare the density profiles with $\varepsilon_0 = 10^{-3}$. Other parameters are chosen to be $\Delta x = 0.1, \Delta t = 10^-2, \chi_N=1.1, \chi_S = 0.5, D_s = 2, \alpha = 0.05$. \textcolor{red}{(The scheme is unstable with a larger $\varepsilon$, say 0.01.) Can we try other values e.g. $2\cdot 10^{-3}, $, $5\cdot 10^{-3}$ to see some differences?}}
%     \label{fig:epsilon}
% \end{figure}
\begin{figure}
    \centering
    \includegraphics[width=0.45\linewidth]{travelling_speed_analytic.eps}
    \includegraphics[width=0.45\linewidth]{travelling_speed_numeric_macro.eps}
    \centering \includegraphics[width=0.45\linewidth]{travelling_speed_numeric_kinetic.eps}
    \caption{Analytical travelling speed with $D_s = 2, \alpha = 0.05$ and different values of $\chi_S$ and $\chi_N$.}
    \label{fig:travelling_speed}
\end{figure}
{{In the simulation, a flux limiter is adopted to make sure that $f$ is nonnegative, which improved the robustness of our scheme. 
Figure \ref{fig:conv_T} compares the density profiles of the kinetic model for $\varepsilon = 0.1, 0.05, 0.01$ with that of the macroscopic model. %\textcolor{red}{(This figure contradicts Fig 3 where you say you cannot choose larger values of $\varepsilon$ and where all solutions look the same, while in this other plot we see some differences.)}
%\XR{(There is no contradiction. In Fig. 3, $\varepsilon$ is much smaller. It seems that we can do numerical computation now with a relatively large $\varepsilon$, as shown in Fig. 4.)}
Table \ref{tab:epsilon_T_values} shows the $L^2$-error between the mesoscopic and macroscopic solutions, $\|\rho^\varepsilon-\rho\|_2$, for different values of $\varepsilon$ and T. 
Table \ref{tab:travel_speed} shows numerical travelling speeds corresponding to various $\varepsilon$ and T. 
}}
\begin{figure}[htbp]
    \centering
    \includegraphics[width=0.4\textwidth]{T_1.eps}
    \includegraphics[width=0.4\textwidth]{T_10.eps}
    \centering
    \includegraphics[width=0.4\textwidth]{T_20.eps}
    \includegraphics[width=0.4\textwidth]{T_50.eps}
    \caption{Comparison of density profiles for $\varepsilon = 0.1$, $0.05$, and $0.01$. The other parameters are set as follows: $\Delta x = 0.1$, $\Delta t = 10^{-3}$, $\chi_N = 1.1$, $\chi_S = 0.5$, $D_s = 2$, and $\alpha = 0.05$.}
    \label{fig:conv_T}
\end{figure}

\begin{table}[h!]
    \centering
    \begin{tabular}{|c|c|c|c|c|c|}
        \hline
        & $\varepsilon= 1$ & $\varepsilon=0.1$ & $\varepsilon=0.05$ & $\varepsilon=0.01$ & $\varepsilon=0.001$ \\
        \hline
        T = 1 & 0.457387598 & 0.195572754 & 0.063367709 & 0.005665137 & 0.004720726 \\
        \hline
        T = 5 & 1.170938306 & 0.166912175 & 0.048322629 & 0.004832707 & 0.002355114 \\
        \hline
        T = 10 & 1.802961097 & 0.143816224 & 0.040773215 & 0.003372185 & 0.00141042 \\
        \hline
        T = 20 & 2.858996388 & 0.143795041 & 0.040367427 & 0.00313379 & 0.000241319\\
        \hline
        T = 50 & 6.209264324 & 0.139749525 & 0.038962193 & 0.003058151 & 0.000246702 \\
        \hline
    \end{tabular}
    \caption{$L^2$-error $\|\rho^\varepsilon - \rho\|_2$ for different values of $\varepsilon$ and $T$. The travelling wave emerges after $T = 20$, and the results in the last two rows are nearly identical.}
    \label{tab:epsilon_T_values}
\end{table}

\begin{table}[h!]
    \centering
    \begin{tabular}{|l|c|c|c|c|c|c|}
        \hline
         & $T=1$ & $T=10$ & $T=20$ & $T=30$ & $T=40$ & $T=50$ \\
        \hline
        Macro & 0.880333028 & 0.599239839 & 0.727553803 & 0.733532824 & 0.725697528 & 0.722771112 \\
        Kinetic $\varepsilon=1$ & 0.518963817 & 0.337444262 & 0.325072197 & 0.328886534 & 0.325139289 & 0.312056651 \\
        Kinetic $\varepsilon=10^{-1}$ & 0.755476966 & 0.622117057 & 0.73996298 & 0.734300279 & 0.72678202 & 0.72959248 \\
        Kinetic $\varepsilon=10^{-2}$ & 0.747591452 & 0.616041174 & 0.729054516 & 0.734959884 & 0.72684429 & 0.724830459 \\
        Kinetic $\varepsilon=10^{-3}$ & 0.816131653 & 0.60631475 & 0.727427547 & 0.733612479 & 0.725641466 & 0.723017182 \\
        \hline
    \end{tabular}
    \caption{Numerical travelling speeds of the macroscopic and kinetic models for various values of $\varepsilon$ and $T$.}
    \label{tab:travel_speed}
\end{table}
\begin{figure}
\centering
\includegraphics[width=0.6\linewidth]{error_vs_time.eps}\includegraphics[width=0.6\linewidth]{wavespeed_vs_time.eps}
    \caption{These plots correspond to Table 3 and 4.}
    \label{fig:placeholder}
\end{figure}

\appendix

\section{Numerical scheme for the PDE}
For the numerical solution of the PDE we rewrite \eqref{eq: rho nondimension} as
\begin{equation}
    \partial_t\rho=\nabla\cdot\Bigl(\rho(A\nabla H'(\rho)-Bu_S-Cu_N)\Bigr)\ .
\end{equation}
\textcolor{red}{(Explain numerical discretisation)}

Then we have
\begin{equation}
\frac{d\rho_j}{dt}=-\frac{F_{j+\frac{1}{2}}(t)-F_{j-\frac{1}{2}}(t)}{\Delta x},
\end{equation}
where
\begin{equation}
F_{j+\frac{1}{2}}=U^+_{j+\frac{1}{2}}\rho_j+U^-_{j+\frac{1}{2}}\rho_{j+1}, \quad  F_{j-\frac{1}{2}}=U^+_{j-\frac{1}{2}}\rho_{j-1}+U^-_{j-\frac{1}{2}}\rho_j\ .
\end{equation}
Also
\[
\xi_i=AH'(\rho)=A\log(\rho_i)
\]
and
\begin{equation}
    U_{j+\frac{1}{2}}=-\frac{\xi_{j+1}-\xi_j}{\Delta x}+a_j[S]+a_j[N]\ ,\ \quad    U_{j-\frac{1}{2}}=-\frac{\xi_{j}-\xi_{j-1}}{\Delta x}+a_{j-1}[S]+a_{j-1}[N]\ ,
\end{equation}
\[
a_j[S]=B\sign\Bigl( \frac{S_{j+1}-S_j}{\Delta x}\Bigr)\ ,\ \quad a_{j-1}[S]=B\sign\Bigl( \frac{S_{j}-S_{j-1}}{\Delta x}\Bigr)
\]
and finally
\begin{equation}
    U_{j+\frac{1}{2}}^+=\max(U_{j+\frac{1}{2}},0)\ ,\quad U_{j+\frac{1}{2}}^-=\min(U_{j+\frac{1}{2}},0)
\end{equation}

We consider zero-flux boundary conditions for $\rho$, $S$ and $N$.



\section{Tumbling operator}\label{app: T}

\begin{figure}[!h]
    \centering
    \includegraphics[width=0.6\linewidth]{tumbling rates.png}
    \caption{In this figure there are plotted the profiles of $S$ and $N$ on the plane defined by the axes $z=0$ and $v=c$. In each quadrant we have the corresponding value of the tumbling rate.}
    \label{tumbling rates}
\end{figure}

To better understand these values, we will explain the intuition behind the signs in each expression of $T$.

\begin{itemize}
    \item $T^+_- = 1-\chi_S-\chi_N$: In this particular case, the $+$ tells us that $v>c$, i.e. the bacteria of this group are moving faster than the wave, while the $-$ corresponds to $z<0$ which means that bacteria are at the left side of the peak of $S$ concentration. Therefore we are in the situation where bacteria move towards the maximum concentration of $S$ and also towards higher concentration of $N$, since it is a favourable motion we do not want to change it that much so we need the tumbling rate as small as possible, so $T^+_-$ is minimum.
    \item $T^+_+ = 1+\chi_S-\chi_N$: Here, we also have $v>c$ which indicates that the bacteria are moving faster than the wave (to the right). As far as the nutrient is concerned, we are going in the right direction, thus, $\chi_N$ does not contribute to more tumbling, on the contrary, wants less tumbling, and that is why we have $-\chi_N$ as well as in the case above. On the other hand, the subscript $+$ shows that $z>0$, i.e. we are in the right side of the peak of $S$. So, as far as the $S$ is concerned, bacteria move away from the higher concentration of $S$ (not favourable for $S$). Consequently, the factor $\chi_S$ does contribute to a higher tumbling rate to revert this situation, that is why we have a positive $+\chi_S$.

    \item $T^-_- = 1+\chi_S+\chi_N$: In the third setup we have $v<c$ which means that we are moving slower or even in the opposite direction to the wave, and $z<0$ indicating bacteria are in the left side of the peak of $S$. We are in the worst scenario because bacteria move away (or goes slowly) towards the peak of $S$ and from the higher concentration of $N$ too. Taking this into account the best is to tumble more to see if we can rectify this situation. Hence, the $T$ is maximal: both factors contribute to have more tumbling.

    \item $T^-_+ = 1-\chi_S+\chi_N$: Finally, we have the setup where the $-\chi_S$ does not contribute to more tumbling because, since bacteria are at the right side of the peak, moving backwards or slower it is fine for $S$. Oppositely, it is not a good situation for $N$, so the factor $+\chi_N$ does contribute to more tumbling.
\end{itemize}

\section{Post-processing for positivity preservation}\label{sec:limiter}
In this appendix, we describe the positivity-preserving post-processing procedure applied at each time step.  
Given the numerical solutions $\tilde{\rho}_{j+\frac12}^{\,n+1}$ and $\tilde{g}_{j,k}^{\,n+1}$ obtained from the scheme~\eqref{eq: main scheme}, the goal is to construct the modified solutions $\rho_{j+\frac12}^{\,n+1}$ and $g_{j,k}^{\,n+1}$ such that the reconstructed distribution
\[
    f_{j+\frac12,k}^{\,n+1} 
    = \rho_{j+\frac12}^{\,n+1}\psi_0(v_k)
    + \varepsilon \bar{g}_{j+\frac12,k}^{\,n+1}
\]
is nonnegative for all $(j,k)$, while the total mass is conserved.

% Step 1: We reconstruct $f$ as 
% $$
%     f_{j,k}^{n+1} = \bar{\rho}_j^{n+1} \psi_0(v_k) + \varepsilon \tilde{g}_{j,k}^{n+1},
% $$
% where $\bar{\rho}_{j}^{n+1} = (\tilde{\rho}_{j+\frac12}^{n+1} + \tilde{\rho}_{j-\frac12}^{n+1}) / 2$. It is easy to see that $<f_{j,k}^{n+1}>_h = \bar{\rho}_j^{n+1}$. 

% Step 2: We modify the solutions so that they are positive. 
% If $\bar{\rho}_j^{n+1}>0$, by a Zhang-Shu type limiter, we set 
% $$
%     \rho_j^{n+1,*} = \bar{\rho}_j^{n+1}, \quad g_{j,k}^{n+1,*}= \theta_j (f_{j,k}^{n+1} - \bar{\rho}_j^{n+1}),
% $$
% where $\theta_j = \min\left\{1, \frac{\bar{\rho}_j^{n+1}}{\bar{\rho}_j^{n+1} - m_j^{n+1}}\right\}$ 
% with $m_j^{n+1} = \min_k \, f_{j,k}^{n+1}$. 
% If $\bar{\rho}_j^{n+1}\le 0$, we force $\rho_j^{n+1,*}=0$ and $g_{j,k}^{n+1,*}=0$.

% Step 3: We make the numerical solutions to be mass conserving without changing the positivity. 
% To make sure that the mass is conserved, we set 
% $$
%    \rho_j^{n+1,**} = \alpha  \rho_j^{n+1,*}, 
%    \quad
%    g_{j,k}^{n+1,**} = \alpha g_{j,k}^{n+1,*},
% $$
% where 
% $$
%     \alpha = \frac{\sum_{j,k} \tilde{f}_{j,k}^{n+1}}{\sum_{j,k} f_{j,k}^{n+1,**}}
% $$
% with $f_{j,k}^{n+1,**} = \rho_j^{n+1,**} \psi_0(v_k) + \varepsilon \tilde{g}_{j,k}^{n+1,**}$. 

% Step 4: The numerical solutions at $t=t_{n+1}$ can be reconstructed as 
% $$
%     \rho_{j+\frac12}^{n+1} = \frac{\rho_j^{n+1,**} + \rho_{j+1}^{n+1,**}}{2}, 
%     \quad
%     g_{j,k}^{n+1} = g_{j,k}^{n+1,**}.
% $$
% In this way, we get the positive numerical solutions without changing mass.


\paragraph{Step 1: Reconstruction of $f$ on cell centers.}
We first define the cell-centered perturbation by
\[
    \bar{g}_{j+\frac12,k}^{\,n+1}
    := \frac{\tilde{g}_{j,k}^{\,n+1}
    + \tilde{\rho}_{j+1,k}^{\,n+1}}{2}.
\]
Then we reconstruct a cell-centered distribution
\[
    \tilde{f}_{j+\frac12,k}^{\,n+1}
    := \tilde{\rho}_{j+\frac12}^{\,n+1}\psi_0(v_k)
     + \varepsilon\,\bar{g}_{j+\frac12,k}^{\,n+1},
\]
which satisfies the consistency relation  
$\langle \tilde{f}_{j+\frac12,\cdot}^{\,n+1}\rangle_h = \tilde{\rho}_{j+\frac12}^{\,n+1}$.

\paragraph{Step 2: Cellwise positivity correction.}
We apply a Zhang--Shu type scaling limiter to remove possible negative undershoots while keeping the cell average unchanged.

Let
\[
    m_{j+\frac12}^{\,n+1} := \min_{k}\, \tilde{f}_{j+\frac12,k}^{\,n+1}.
\]
If $\tilde{\rho}_{j+\frac12}^{\,n+1} > 0$, we set
\[
    \rho_{j+\frac12}^{\,n+1,*} = \tilde{\rho}_{j+\frac12}^{\,n+1}, \qquad
    f_{j+\frac12,k}^{\,n+1,*}
    = \tilde{\rho}_{j+\frac12}^{\,n+1} + \theta_{j+\frac12} \big(\tilde{f}_{j+\frac12,k}^{\,n+1} - \tilde{\rho}_{j+\frac12}^{\,n+1}\big),
\]
with scaling factor
\[
    \theta_{j+\frac12}
    := \min\left\{
        1,\,
        \frac{\tilde{\rho}_{j+\frac12}^{\,n+1}}
        {\tilde{\rho}_{j+\frac12}^{\,n+1} - m_{j+\frac12}^{\,n+1}}
    \right\}.
\]
This ensures $f_{j+\frac12,k}^{\,n+1,*} \ge 0$ for all $k$, while preserving the cell average.

If $\tilde{\rho}_{j+\frac12}^{\,n+1}\le 0$, positivity cannot be recovered without changing the cell average.  
In this case, we set
\[
    \rho_{j+\frac12}^{\,n+1,*}=0,\qquad
    g_{j+\frac12,k}^{\,n+1,*}=0,
\]
which corresponds to removing the (nonphysical) negative mass in that cell.

\paragraph{Step 3: Global mass correction.}
After Step~2, the modified distribution
\begin{equation}\label{eq:modified_f}
    f_{j+\frac12,k}^{\,n+1,*}
    = \rho_{j+\frac12}^{\,n+1,*}\psi_0(v_k)
    + \varepsilon g_{j+\frac12,k}^{\,n+1,*}
\end{equation}
is nonnegative but may not conserve the total mass.  
To enforce mass conservation, we apply a uniform scaling:
\[
    \rho_{j+\frac12}^{\,n+1,**} = \alpha\,\rho_{j+\frac12}^{\,n+1,*},
    \qquad
    g_{j+\frac12,k}^{\,n+1,**} = \alpha\,g_{j+\frac12,k}^{\,n+1,*},
\]
where the factor $\alpha$ is chosen as
\[
    \alpha
    := \frac{\displaystyle\sum_{j}\tilde{\rho}_{j+\frac12}^{\,n+1}\,\Delta x}
            {\displaystyle\sum_{j}\rho_{j+\frac12}^{\,n+1,*}\,\Delta x},
\]
so that the global mass before and after limiting is identical.

\paragraph{Step 4: Reconstruction on staggered grids.}
Finally, the staggered density and microscopic variables used by the main scheme are reconstructed as
\[
    \rho_{j+\frac12}^{\,n+1}
    := \rho_{j+\frac12}^{\,n+1,**},
    \qquad
    g_{j,k}^{\,n+1} := \frac{g_{j+\frac12,k}^{\,n+1,**}+g_{j-\frac12,k}^{\,n+1,**}}{2}.
\]
where $g_{j+\frac12,k}^{\,n+1,**}$ can be solved from \eqref{eq:modified_f}. 

The resulting pair $(\rho_{j+\frac12}^{\,n+1}, g_{j,k}^{\,n+1})$ satisfies:
(i) $f_{j,k}^{\,n+1}\ge 0$ for all $(j,k)$;  
(ii) the total mass is conserved;  
(iii) $\langle g_{j,\cdot}^{\,n+1}\rangle_h = 0$ for all $j$.

\iffalse
IN THE EXAMPLE\\
$\epsilon_j=D_p \log \rho_j-\chi(S) sign(\frac{\partial S}{\partial x})X-\chi(N) sign(\frac{\partial N}{\partial x})X$
\begin{eqnarray}
U^+_{j+\frac{1}{2}}=\max(U_{j+\frac{1}{2}},0), \quad U^-_{j+\frac{1}{2}}=\min (U_{j+\frac{1}{2}},0)
\end{eqnarray}

\begin{eqnarray}
\frac{S_j^{n+1}-S_j^n}{\Delta t}=D_S\frac{M_{j+1}^n-2S_j^n+S_{j-1}^n}{(\Delta x)^2}-\alpha S_j^n+\beta \rho_j^n\\
S_j^{n+1}=S_j^n+D_S \frac{\Delta t} {(\Delta x)^2}(S_{j+1}-2S_j^n+S_{j-1}^n)-\alpha \Delta t S_j^n+\beta \Delta t \rho_j^n
\end{eqnarray}


\begin{eqnarray}
\frac{N_j^{n+1}-N_j^n}{\Delta t}=D_N\frac{N_{j+1}^n-2N_j^n+N_{j-1}^n}{(\Delta x)^2}-\gamma N_j^n \rho_j^n\\
N_j^{n+1}=N_j^n+D_N \frac{\Delta t} {(\Delta x)^2}(N_{j+1}-2N_j^n+N_{j-1}^n)-\gamma \Delta t N_j^n \rho_j^n
\end{eqnarray}

\textbf{ For boundary $ j=1(j=1,...,m) \quad b=D_S \frac{\Delta t} {(\Delta x)^2}, \quad d=D_N \frac{\Delta t} {(\Delta x)^2}$}
\begin{eqnarray}
\frac{d\rho_1}{dt}=-\frac{F_{1+\frac{1}{2}}}{\Delta x}=-\frac{U^+_{\frac{3}{2}}p_1+U^-_{\frac{3}{2}}p_2 }{\Delta x}, \quad U_{\frac{3}{2}}=-\frac{\epsilon_2-\epsilon_1}{\Delta x}
\end{eqnarray}

 
%FOR S
\begin{eqnarray}
\partial_x S |_{x=0}= 0, \quad -> S^n_1-S^n_0=0\\
S_j^{n+1}=S_j^n+D_M \frac{\Delta t} {(\Delta x)^2}(S^n_{j+1}-2M_j^n+S_{j-1}^n)-\alpha \Delta t S_j^n+\beta \Delta t \rho_j^n, \\
S_1^{n+1}=S_1^n+b(S^n_{2}-2S_1^n+S_{0}^n)-\alpha \Delta t S_1^n+\beta \Delta t \rho_1^n\\
S_1^{n+1}=S_1^n+b(S^n_{2}-S_1^n)-\alpha \Delta t S_1^n+\beta \Delta t \rho_1^n\\
S_1^{n+1}=bS^n_2+(1-b-\alpha \Delta t)S_1^n+\beta \Delta t \rho_1^n
\end{eqnarray}

%FOR N
\begin{eqnarray}
\partial_x N |_{x=0}= 0, \quad -> N^n_1-N^n_0=0\\
N_j^{n+1}=N_j^n+D_N \frac{\Delta t} {(\Delta x)^2}(N^n_{j+1}-2N_j^n+N_{j-1}^n)-\gamma \Delta t N_j^n \rho_j^n\\
N_1^{n+1}=N_1^n+d(N^n_{2}-2N_1^n+N_{0}^n)-\gamma \Delta t N_1^n \rho_1^n\\
N_1^{n+1}=N_1^n+d(N^n_{2}-N_1^n)-\gamma \Delta t N_1^n \rho_1^n\\
N_1^{n+1}=dN^n_{2}+(1-d)N_1^n-\gamma \Delta t N_1^n\rho_1^n\\
\end{eqnarray}


\textbf{ For boundary $ j=m(j=1,...,m) \quad b=D_S \frac{\Delta t} {(\Delta x)^2}, \quad d=D_N \frac{\Delta t} {(\Delta x)^2}$}
\begin{eqnarray}
\frac{d\rho_m}{dt}=\frac{F_{m-\frac{1}{2}}}{\Delta x}=\frac{U^+_{m-\frac{1}{2} }\rho_{m-1}+U^-_{m-\frac{1}{2} }\rho_m }{\Delta x}, \quad U_{m-\frac{1}{2}}=-\frac{\epsilon_m-\epsilon_{m-1}}{\Delta x}
\end{eqnarray}


%FOR S
\begin{eqnarray}
\partial_x S |_{x=L}= 0, \quad -> S^n_{m+1}-S^n_m=0\\
S_j^{n+1}=S_j^n+D_S \frac{\Delta t} {(\Delta x)^2}(S^n_{j+1}-2S_j^n+S_{j-1}^n)-\alpha \Delta t S_j^n+\beta \Delta t \rho_j^n, \\
S_m^{n+1}=S_m^n+b(S^n_{m+1}-2S_m^n+S_{m-1}^n)-\alpha \Delta t M_m^n+\beta \Delta t \rho_m^n\\
S_m^{n+1}=S_m^n+b(-S^n_{m}+S_{m-1}^n)-\alpha \Delta t S_m^n+\beta \Delta t \rho_m^n\\
S_m^{n+1}=(1-b-\alpha \Delta t)S^n_m+bS_{m-1}^n+\beta \Delta t \rho_m^n
\end{eqnarray}

%FOR N
\begin{eqnarray}
N(t,x=L)= \tilde{N}, \\
N_j^{n+1}=N_j^n+d(N^n_{j}-2N_{j-1}^n+N_{j-2}^n)-\gamma \Delta t N_j^n \rho_j^n\\
N_m^{n+1}=N_m^n+d(N^n_{m}-2N_{m-1}^n+N_{m-2}^n)-\gamma \Delta t N_m^n \rho_m^n\\
N_m^{n+1}=(1+d-\gamma \Delta t\rho_m^n)\tilde{N}+d(-2N^n_{m-1}+N_{m-2}^n)
\end{eqnarray}


\textbf{\textit{Example}} J. Saragosti, v. Calvez, N. Bournaveas, A. Buguin, P. Silberzan, B. Perthame ''Mathematical Description of Bacterial Travelling Pulses'' 2009\\
$D_{\rho}=1, \chi_M=\chi_N=1$,  $D_M=2,\alpha=0.05,\gamma=\beta=1 $  $D_N=0$ \\
$\Delta x=0.5,\quad \Delta t=0.001,\quad T=5, \quad x\in[0,50]$\\

$\rho(0,x)=\exp(-x),\quad S(0,x)=0,\quad N(0,x)=N_0=10$\\
\fi

\iffalse
SCALING OF PDE\\

$t=\tilde{t} \bar{t}, \quad x=\tilde{x}\bar{x},\quad  \tilde{t},\quad \tilde{x}$ are new variables.

$$
\tilde{t}=\frac{t}{\bar{t}}, \quad \tilde{x}=\frac{x}{\bar{x}}
$$

\begin{eqnarray}
\frac{\partial \rho}{\partial t}&=&D_p\frac{\partial^2 \rho}{\partial x^2}-\frac{\partial}{\partial x}\left( \rho \chi_S sign(\frac{\partial S}{\partial x})+\rho \chi_N sign(\frac{\partial N}{\partial x})\right) \\
\frac{\partial S}{\partial t}&=&D_S \frac{\partial^2S}{\partial x^2}-\alpha S+\beta \rho\\
\frac{\partial N}{\partial t}&=&D_N \frac{\partial^2N}{\partial x^2}-\gamma \rho N
\end{eqnarray}
%FOR S
\begin{eqnarray*}
\frac{1}{\bar{t}}\frac{\partial S}{\partial\tilde{t}}=\frac{D_S}{(\bar{x})^2}\frac{\partial^2S}{\partial \tilde{x}^2}-\alpha S+\beta \rho
\end{eqnarray*}

\begin{eqnarray}
\frac{\partial S}{\partial\tilde{t}}=\frac{D_S\bar{t}}{(\bar{x})^2}\frac{\partial^2S}{\partial \tilde{x}^2}-\alpha\bar{t} S+\beta \bar{t}\rho
\end{eqnarray}

%FOR N

\begin{eqnarray*}
	\frac{1}{\bar{t}}\frac{\partial N}{\partial\tilde{t}}=\frac{D_N}{(\bar{x})^2}\frac{\partial^2N}{\partial \tilde{x}^2}-\gamma\rho N
\end{eqnarray*}

\begin{eqnarray}
\frac{\partial N}{\partial\tilde{t}}=\frac{D_N \bar{t}}{(\bar{x})^2}\frac{\partial^2N}{\partial \tilde{x}^2}-\gamma\bar{t}\rho N
\end{eqnarray}




\begin{eqnarray*}
\frac{1}{\bar{t}}\frac{\partial \rho}{\partial\tilde{t}}=\frac{D_\rho}{(\bar{x})^2}\frac{\partial^2\rho}{\partial \tilde{x}^2}-\frac{1}{\bar{x}}\frac{\partial}{\partial \tilde{x}}\left[ \rho\chi_S sign\left( \frac{\partial S}{\bar{x}\partial \tilde{x}}\right) +\rho \chi_N sign\left( \frac{\partial N}{\bar{x}\partial \tilde{x}}\right)\right] 
\end{eqnarray*}

\begin{eqnarray}
\frac{\partial \rho}{\partial\tilde{t}}=\frac{\bar{t}D_\rho}{(\bar{x})^2}\frac{\partial^2\rho}{\partial \tilde{x}^2}-\frac{1}{\bar{x}}\frac{\partial}{\partial \tilde{x}}\left[ \rho\frac{\bar{t}}{\bar{x}}\chi_S sign\left( \frac{\partial S}{\bar{x}\partial \tilde{x}}\right) +\rho\frac{\bar{t}}{\bar{x}} \chi_N sign\left( \frac{\partial N}{\bar{x}\partial \tilde{x}}\right)\right] 
\end{eqnarray}
FOR $\rho$
\begin{eqnarray}
A=\frac{\bar{t}D_\rho}{(\bar{x})^2},\quad
B= \frac{\bar{t}}{\bar{x}}\chi_S,\quad
 C=\frac{\bar{t}}{\bar{x}} \chi_N
 \end{eqnarray}
 FOR $S$
 \begin{eqnarray}
 D=\frac{\bar{t}D_S}{(\bar{x})^2},\quad
 E=\bar{t}\alpha,\quad
  F=\bar{t}\beta,
\end{eqnarray}
FOR $N$
\begin{eqnarray}
 G=\frac{\bar{t}D_N}{(\bar{x})^2},\quad
 H=\gamma \bar{t}
\end{eqnarray}

\begin{eqnarray*}
A=\frac{\bar{t}D_\rho}{(\bar{x})^2}=\frac{100s}{(200.10^{-4}cm)^2}4.10^{-6}\frac{cm^2}{s}=1,\\
B= \frac{\bar{t}}{\bar{x}}\chi_S=\frac{100s}{200.10^{-4}cm}10^{-4}cms^{-1}=\frac{1}{2}\\,\end{eqnarray*}
\begin{eqnarray*}
C=\frac{\bar{t}}{\bar{x}} \chi_N=\frac{100s}{200.10^{-4}cm}.2.2.10^{-4}cms^{-1}=1.1,\\
D=\frac{\bar{t}D_S}{(\bar{x})^2}=\frac{8.10^{-6}cm^2s^{-1}.100 s}{(200.10^{-4}cm)^2}=2,\\
E=\bar{t}\alpha=100s.4.10^{-3}s^{-1}=0.4\\
F=\bar{t}\beta=100s.4.10^5cell^{-1}s^{-1}=4.10^7cell^{-1}\\
G=\frac{\bar{t}D_N}{(\bar{x})^2}=8.10^{-6}cm^2s^{-1}\frac{100s}{(200.10^{-4}cm)^2}=2,\\
H=\gamma\bar{t}=4.10^{-5}cell^{-1}s^{-1}100s=4.10^{7}cell^{-1}
\end{eqnarray*}

% \begin{table}[]
%     \centering
%     \begin{tabular}{c|c|c|c}
%         \hline
%         $\varepsilon$ & $10^{-3}$ & $10^{-4}$ & $10^{-5}$ \\
%         \hline
%         $\|\rho^{\varepsilon} - \rho\|_2$ & 6.7634e-04 & 6.7943e-05 & 6.7975e-06 \\
%         \hline
%         $\|\rho^{\varepsilon} - \rho\|_{\infty}$ & 
%         0.0013 & 1.3498e-04 & 1.3504e-05 \\
%         \hline
%     \end{tabular}
%     \caption{Accuracy test at $T = 0.1$. \textcolor{red}{(We can't observe the convergence as $\varepsilon\to0$ with a larger $T$. The reason is unclear but might be related to the boundary condition as oscillations appear near $x=0$, which is different from the macro model.)}}
%     \label{tab:accuracy_test}
% \end{table}
\fi

\bibliographystyle{abbrv}
\bibliography{bib.bib}


\end{document}